\documentclass[12pt,a4paper]{article}
\usepackage[T1]{fontenc}
\usepackage[utf8]{inputenc}
\usepackage{lmodern}
\usepackage{amsmath,amssymb,amsthm,mathtools}
\usepackage{geometry}
\geometry{margin=1in}
\usepackage{microtype}
\usepackage{hyperref}
\usepackage{xcolor}
\usepackage{enumitem}
\usepackage{fancyhdr}
\usepackage{bookmark}

\usepackage{centernot}

\DeclareMathOperator{\rank}{rank}

\DeclareMathOperator{\Deck}{Deck}

\usepackage{tikz-cd}

\usepackage{tocloft}
\setlength{\cftsubsecnumwidth}{3.2em}

\newtheorem{prop}{Proposition}[section]

\newcommand{\Q}{\mathbb{Q}}
\newcommand{\RP}{\mathbb{R}\mathrm{P}}

\newcommand{\Z}{\mathbb{Z}}

\newcommand{\Chi}{\chi}
\newcommand{\del}{\partial}

\newtheorem{definition}{Definition}[section]

\newtheorem{example}[definition]{Example}

\hypersetup{
	colorlinks=true,
	linkcolor=blue!50!black,
	urlcolor=blue!50!black,
	citecolor=blue!50!black,
	pdftitle={Theory of Algebraic Topology 8},
	pdfauthor={},
	pdfsubject={Solved problems and theory notes in commutative algebra},
	pdfcreator={OpenAI}
}

\setlist[itemize]{topsep=4pt,itemsep=2pt,parsep=0pt,partopsep=0pt}
\setlength{\parskip}{0.55em}
\setlength{\parindent}{0pt}

\setcounter{secnumdepth}{2}


\setcounter{tocdepth}{2}

\pagestyle{fancy}


\fancyhf{}
\lhead{Theory of Algebraic Topology 8}
\rhead{\thepage}
\renewcommand{\headrulewidth}{0.4pt}

\newtheoremstyle{notespace}{6pt}{6pt}{\normalfont}{} {\bfseries}{.}{0.5em}{}
\theoremstyle{notespace}

\newtheorem{remark}{Remark}[section]

\DeclareMathOperator{\Spec}{Spec}
\DeclareMathOperator{\Ass}{Ass}
\DeclareMathOperator{\Ann}{Ann}
\DeclareMathOperator{\Supp}{Supp}
\DeclareMathOperator{\Hom}{Hom}
\DeclareMathOperator{\Min}{Min}
\DeclareMathOperator{\Ker}{Ker}
\DeclareMathOperator{\Frac}{Frac}

\DeclareMathOperator{\depth}{depth}
\DeclareMathOperator{\Tor}{Tor}
\DeclareMathOperator{\Ext}{Ext}

\DeclareMathOperator{\coker}{coker}
\DeclareMathOperator{\im}{im}

\newcommand{\longdash}{\mathrel{\text{---}}}

\title{\vspace{-1.5em}\bfseries Theory of Algebraic Topology 8}
\author{}
\date{\today}

\begin{document}
	\maketitle
	\thispagestyle{plain}
		\begin{center}
		\begin{minipage}{0.92\textwidth}
			Lecture notes with worked examples in algebraic topology
		\end{minipage}
	\end{center}
		\pdfbookmark[1]{Contents}{contents}\tableofcontents
	\clearpage

\section{Euler characteristic: definition, intuition, and examples}

The Euler characteristic is one of the simplest and most useful numerical invariants in topology. It is easy to compute from a triangulation or cell decomposition, and at the same time it detects surprisingly deep information about the shape of a space.

\subsection{First intuition: vertices, edges, and faces}

For a polygonal or polyhedral object, the Euler characteristic is first seen through the alternating count
\[
\chi = V-E+F,
\]
where
\[
V=\text{number of vertices},\qquad
E=\text{number of edges},\qquad
F=\text{number of faces}.
\]

For a planar polygonal region or a polyhedron, this number often stays unchanged when one subdivides the figure into smaller pieces. That is the first sign that \(\chi\) depends on the underlying shape, not on the particular way we cut it up.

\subsection{Definition from a triangulation}

\begin{definition}
	Let \(X\) be the polyhedron of a finite simplicial complex \(K\). If \(c_i\) denotes the number of \(i\)-simplices of \(K\), then the \emph{Euler characteristic} of \(X\) is
	\[
	\chi(X)=\sum_{i\ge 0}(-1)^i c_i.
	\]
	Thus
	\[
	\chi(X)=c_0-c_1+c_2-c_3+\cdots.
	\]
\end{definition}

In low dimensions this becomes:
\[
\chi(X)=V-E+F
\]
for a triangulated surface or planar complex, and
\[
\chi(X)=V-E+F-T
\]
for a triangulated \(3\)-dimensional complex, where \(T\) is the number of tetrahedra.

\subsection{Definition from a cell decomposition}

More generally, if a space \(X\) has a finite CW decomposition with \(c_i\) cells in dimension \(i\), then
\[
\chi(X)=\sum_{i\ge 0}(-1)^i c_i.
\]

This is very useful because many spaces are easier to describe by cells than by simplices. For example, a torus can be built from one \(0\)-cell, two \(1\)-cells, and one \(2\)-cell, so
\[
\chi(T^2)=1-2+1=0.
\]

\subsection{Why the Euler characteristic is topological}

A fundamental fact is that \(\chi(X)\) does not depend on the particular triangulation or cell decomposition. Different triangulations of the same space always give the same Euler characteristic.

So one may compute \(\chi(X)\) from whichever decomposition is easiest.

\subsection{First examples: polygons and planar regions}

\subsubsection*{Example 1: a line segment}

A line segment has two vertices and one edge, so
\[
\chi([0,1])=2-1=1.
\]

If we subdivide it by adding one extra vertex in the middle, then we get three vertices and two edges:
\[
\chi=3-2=1.
\]
So subdivision does not change the value.

\subsubsection*{Example 2: a triangle together with its interior}

A filled triangle has
\[
V=3,\qquad E=3,\qquad F=1.
\]
Hence
\[
\chi(\Delta^2)=3-3+1=1.
\]

This is the same value as for a disk, since a filled triangle is topologically a disk.

\subsubsection*{Example 3: a square together with its interior}

A square has
\[
V=4,\qquad E=4,\qquad F=1,
\]
so
\[
\chi(\text{square})=4-4+1=1.
\]

Again, this matches the fact that a square is topologically a disk.

\subsubsection*{Example 4: a pentagon together with its interior}

A filled pentagon has
\[
V=5,\qquad E=5,\qquad F=1,
\]
so
\[
\chi=5-5+1=1.
\]

Thus every filled \(n\)-gon has Euler characteristic
\[
\chi=1.
\]

This is a very good first family of examples: all filled polygons are topologically disks.

\subsubsection*{Example 5: the boundary of a polygon}

Now take only the boundary of an \(n\)-gon. Then
\[
V=n,\qquad E=n,\qquad F=0,
\]
hence
\[
\chi=n-n=0.
\]

So the boundary of any polygon has Euler characteristic \(0\). This matches the fact that the boundary is topologically a circle:
\[
\chi(S^1)=0.
\]

\subsection{Examples with triangulations}

\subsubsection*{Example 6: a disk triangulated into two triangles}

Take a square and draw one diagonal. Then the square is decomposed into two triangles. The counts are
\[
V=4,\qquad E=5,\qquad F=2.
\]
Hence
\[
\chi=4-5+2=1.
\]

So even after triangulation, a disk still has Euler characteristic \(1\).

\subsubsection*{Example 7: a sphere as the boundary of a tetrahedron}

The boundary of a tetrahedron is a triangulation of \(S^2\). It has
\[
V=4,\qquad E=6,\qquad F=4.
\]
Therefore
\[
\chi(S^2)=4-6+4=2.
\]

This is one of the most important basic values in topology:
\[
\chi(S^2)=2.
\]

\subsubsection*{Example 8: a torus from a square model}

A torus can be obtained from a square by identifying opposite sides:
\[
T^2=[0,1]^2/(x,0)\sim(x,1),\ (0,y)\sim(1,y).
\]

In the quotient, all four vertices become one point, the horizontal edges become one \(1\)-cell, the vertical edges become one \(1\)-cell, and the square interior becomes one \(2\)-cell. Thus
\[
c_0=1,\qquad c_1=2,\qquad c_2=1.
\]
So
\[
\chi(T^2)=1-2+1=0.
\]

\subsubsection*{Example 9: the projective plane from a square model}

The projective plane can be obtained from a square with opposite sides identified in the pattern
\[
aba b
\]
up to orientation convention. In a standard CW decomposition one gets
\[
c_0=1,\qquad c_1=1,\qquad c_2=1.
\]
Hence
\[
\chi(\RP^2)=1-1+1=1.
\]

\subsubsection*{Example 10: the Klein bottle}

A standard CW decomposition of the Klein bottle has
\[
c_0=1,\qquad c_1=2,\qquad c_2=1.
\]
Thus
\[
\chi(K)=1-2+1=0.
\]

So the torus and the Klein bottle both have Euler characteristic \(0\), even though one is orientable and the other is not.

\subsection{Euler characteristic of standard surfaces}

For closed orientable surfaces \(\Sigma_g\),
\[
\chi(\Sigma_g)=2-2g.
\]

So:
\[
\chi(S^2)=\chi(\Sigma_0)=2,
\]
\[
\chi(T^2)=\chi(\Sigma_1)=0,
\]
\[
\chi(\Sigma_2)=-2,
\]
and in general each additional handle decreases \(\chi\) by \(2\).

For closed nonorientable surfaces \(S_n\),
\[
\chi(S_n)=2-n.
\]

So:
\[
\chi(\RP^2)=\chi(S_1)=1,
\]
\[
\chi(\text{Klein bottle})=\chi(S_2)=0,
\]
\[
\chi(S_3)=-1,
\]
and each additional crosscap decreases \(\chi\) by \(1\).

\subsection{Why polygons are a good starting point}

Polygons and their boundaries are the best first examples because they already show the two most basic values:
\[
\chi(D^2)=1,\qquad \chi(S^1)=0.
\]

More precisely:
\begin{itemize}
	\item every filled polygon is topologically a disk, so it has Euler characteristic \(1\),
	\item every polygon boundary is topologically a circle, so it has Euler characteristic \(0\).
\end{itemize}

From there, one passes naturally to:
\begin{itemize}
	\item polyhedra such as the tetrahedron boundary for the sphere,
	\item edge identifications on polygons for the torus, Klein bottle, and projective plane,
	\item triangulated surfaces of higher genus.
\end{itemize}

So polygons are the simplest entry point into the full theory.

\subsection{Basic properties of Euler characteristic}

\subsubsection*{Property 1: invariance under subdivision}

Subdividing a triangulation does not change \(\chi\). For example, subdividing one edge adds one vertex and one edge, so the quantity
\[
V-E+F
\]
does not change.

\subsubsection*{Property 2: disjoint unions}

If \(X\) and \(Y\) are finite complexes, then
\[
\chi(X\sqcup Y)=\chi(X)+\chi(Y).
\]

For example, two disjoint circles have Euler characteristic
\[
0+0=0.
\]

\subsubsection*{Property 3: inclusion--exclusion}

If \(X=A\cup B\) and all spaces involved have well-defined Euler characteristic, then
\[
\chi(X)=\chi(A)+\chi(B)-\chi(A\cap B).
\]

For example, if two disks overlap in a disk, then
\[
\chi(A\cup B)=1+1-1=1,
\]
which is again the Euler characteristic of a disk.

\subsubsection*{Property 4: connected sums of surfaces}

If \(X\) and \(Y\) are closed surfaces, then
\[
\chi(X\#Y)=\chi(X)+\chi(Y)-2.
\]

This is because one removes a disk from each surface and glues along the resulting boundary circles.

For example,
\[
\chi(T^2\#\RP^2)=\chi(T^2)+\chi(\RP^2)-2=0+1-2=-1,
\]
which matches
\[
\chi(S_3)=2-3=-1.
\]

\subsubsection*{Property 5: finite coverings}

If
\[
p:\widetilde X\to X
\]
is an \(n\)-sheeted covering of finite triangulated spaces, then
\[
\chi(\widetilde X)=n\,\chi(X).
\]

For surfaces this gives strong restrictions. For example, if
\[
\Sigma_g\to \Sigma_h
\]
is an \(n\)-sheeted covering, then
\[
2-2g=n(2-2h).
\]

\subsection{Homological interpretation}

There is also a homological formula:
\[
\chi(X)=\sum_{i\ge 0}(-1)^i \rank H_i(X;\Q).
\]

This is often called the Euler--Poincar\'e formula. It says that the alternating count of simplices equals the alternating count of Betti numbers.

For example, for the torus:
\[
H_0(T^2;\Q)\cong \Q,\qquad
H_1(T^2;\Q)\cong \Q^2,\qquad
H_2(T^2;\Q)\cong \Q,
\]
so
\[
\chi(T^2)=1-2+1=0.
\]

For the sphere:
\[
H_0(S^2;\Q)\cong \Q,\qquad
H_1(S^2;\Q)=0,\qquad
H_2(S^2;\Q)\cong \Q,
\]
so
\[
\chi(S^2)=1+1=2.
\]

\subsection{A compact table of standard values}

\[
\begin{array}{c|c}
	X & \chi(X)\\ \hline
	\text{point} & 1\\
	\left[0,1\right] & 1\\
	D^2 & 1\\
	S^1 & 0\\
	S^2 & 2\\
	T^2 & 0\\
	\RP^2 & 1\\
	\text{Klein bottle} & 0\\
	\Sigma_g & 2-2g\\
	S_n & 2-n
\end{array}
\]

\subsection{Summary}

The Euler characteristic begins as the simple alternating count
\[
V-E+F
\]
for polygons and polyhedra, but it extends to all finite simplicial complexes and CW complexes by
\[
\chi(X)=\sum_i (-1)^i c_i.
\]

The best first examples are:
\[
\chi(\text{filled polygon})=1,
\qquad
\chi(\text{polygon boundary})=0,
\]
followed by
\[
\chi(S^2)=2,\qquad \chi(T^2)=0,\qquad \chi(\RP^2)=1.
\]

These examples already contain the main geometric patterns:
\begin{itemize}
	\item contractible spaces usually have Euler characteristic \(1\),
	\item circles have Euler characteristic \(0\),
	\item adding a handle lowers \(\chi\) by \(2\),
	\item adding a crosscap lowers \(\chi\) by \(1\).
\end{itemize}


\section{Polygonal models of surfaces and how Euler characteristic is read from edge identifications}

A very efficient way to describe many closed surfaces is to start with a polygon and identify its boundary edges in pairs. This gives a concrete geometric model for the surface and makes the Euler characteristic especially easy to compute.

\subsection{Basic idea of a polygonal model}

Start with a closed polygon \(P\), usually a square or a \(2m\)-gon. The interior of \(P\) will become a single \(2\)-cell. The edges of the boundary are then identified in pairs according to a prescribed pattern.

After the identifications:
\begin{itemize}
	\item the whole interior of the polygon gives one \(2\)-cell,
	\item each pair of identified boundary edges gives one \(1\)-cell in the quotient,
	\item boundary vertices may collapse together, producing one or more \(0\)-cells.
\end{itemize}

Thus the Euler characteristic can often be read directly from:
\[
\chi = c_0-c_1+c_2.
\]

Since the polygon interior always contributes one \(2\)-cell, we usually have
\[
c_2=1.
\]

\subsection{How edge labels are written}

To describe identifications, one labels the boundary edges of the polygon cyclically. For example:
\[
aba^{-1}b^{-1}
\]
means that the first and third edges are identified with opposite orientation, and the second and fourth edges are identified with opposite orientation.

Likewise:
\[
aba b
\]
means that the two \(a\)-edges are glued in the same direction, and the two \(b\)-edges are also glued in the same direction.

The orientation matters: opposite orientations usually correspond to orientable gluing, while same-direction gluing introduces nonorientability.

\subsection{First warm-up: a cylinder from a rectangle}

Take a rectangle and identify only the left and right sides in the same vertical direction. The top and bottom remain unglued.

Then:
\begin{itemize}
	\item the interior is one \(2\)-cell,
	\item the left and right edges become one \(1\)-cell,
	\item the top and bottom edges remain distinct boundary edges,
	\item vertices are identified in pairs.
\end{itemize}

The resulting space is a cylinder, which is not closed because it still has boundary.

This is a good first example because it shows that not every edge-identification produces a closed surface. To get a closed surface, every boundary edge must be glued to another boundary edge.

\subsection{The torus from a square}

Take a square and label its edges
\[
aba^{-1}b^{-1}.
\]
This means:
\begin{itemize}
	\item the top and bottom edges are glued with opposite directions,
	\item the left and right edges are glued with opposite directions.
\end{itemize}

The quotient is the torus \(T^2\).

\subsubsection*{Cell count}

After gluing:
\begin{itemize}
	\item all four vertices become one vertex,
	\item the two horizontal edges become one \(1\)-cell,
	\item the two vertical edges become one \(1\)-cell,
	\item the square interior becomes one \(2\)-cell.
\end{itemize}

So:
\[
c_0=1,\qquad c_1=2,\qquad c_2=1.
\]
Hence
\[
\chi(T^2)=1-2+1=0.
\]

\subsubsection*{Why this matches the general formula}

Since \(T^2=\Sigma_1\), the orientable surface formula gives
\[
\chi(\Sigma_1)=2-2\cdot 1=0,
\]
which agrees perfectly.

\subsection{The projective plane from a square}

A standard polygonal model for the projective plane is obtained from a square with word
\[
aa.
\]
Equivalently, one may imagine a \(2\)-gon with its two boundary edges identified in the same direction. Using a square picture is often more convenient visually.

The key feature is that the two copies of the boundary edge are glued with the same orientation, which produces nonorientability.

\subsubsection*{Cell count}

In the standard CW decomposition:
\begin{itemize}
	\item all vertices become one \(0\)-cell,
	\item the identified boundary gives one \(1\)-cell,
	\item the interior gives one \(2\)-cell.
\end{itemize}

Thus:
\[
c_0=1,\qquad c_1=1,\qquad c_2=1.
\]
So
\[
\chi(\RP^2)=1-1+1=1.
\]

\subsubsection*{Comparison with the nonorientable formula}

Since \(\RP^2=S_1\), we also have
\[
\chi(S_1)=2-1=1.
\]

\subsection{The Klein bottle from a square}

Take a square and identify edges according to the word
\[
aba b.
\]

One pair is glued in the orientable way, and the other pair in the nonorientable way. The result is the Klein bottle.

\subsubsection*{Cell count}

After identification:
\begin{itemize}
	\item all four corners become one vertex,
	\item the two \(a\)-edges become one \(1\)-cell,
	\item the two \(b\)-edges become one \(1\)-cell,
	\item the square interior becomes one \(2\)-cell.
\end{itemize}

So:
\[
c_0=1,\qquad c_1=2,\qquad c_2=1.
\]
Hence
\[
\chi(K)=1-2+1=0.
\]

\subsubsection*{Comparison with surface notation}

The Klein bottle is \(S_2\), so the nonorientable formula gives
\[
\chi(S_2)=2-2=0.
\]
Again this agrees.

\subsection{Why the torus and Klein bottle have the same Euler characteristic}

Both spaces admit polygonal models with:
\[
c_0=1,\qquad c_1=2,\qquad c_2=1.
\]
So both have Euler characteristic \(0\).

However, they are not homeomorphic, because:
\begin{itemize}
	\item the torus is orientable,
	\item the Klein bottle is nonorientable.
\end{itemize}

This is an important lesson: Euler characteristic is useful, but it does not classify spaces by itself.

\subsection{The sphere from a polygonal point of view}

Although the sphere is usually pictured as the boundary of a tetrahedron or other polyhedron, it also fits naturally into the polygonal philosophy.

A sphere can be obtained by gluing two disks along their boundary circle. In terms of CW structure, one may take:
\begin{itemize}
	\item one \(0\)-cell,
	\item no \(1\)-cells,
	\item one \(2\)-cell attached trivially above and below, or equivalently a decomposition with
	\[
	c_0=1,\quad c_1=0,\quad c_2=1,\quad \text{plus another \(2\)-cell in a two-disk picture.}
	\]
\end{itemize}

The more standard triangulated model is the boundary of a tetrahedron:
\[
V=4,\qquad E=6,\qquad F=4,
\]
so
\[
\chi(S^2)=4-6+4=2.
\]

Thus the sphere has Euler characteristic \(2\), larger than all other standard closed surfaces.

\subsection{Orientable surface of genus \(g\) from a \(4g\)-gon}

A standard polygonal model for the closed orientable surface \(\Sigma_g\) is:
\[
a_1b_1a_1^{-1}b_1^{-1}\cdots a_gb_ga_g^{-1}b_g^{-1}.
\]

This is a \(4g\)-gon.

\subsubsection*{Cell count}

In the quotient:
\begin{itemize}
	\item all vertices identify to one vertex,
	\item each pair of matching edges gives one \(1\)-cell, so there are \(2g\) such cells,
	\item the interior gives one \(2\)-cell.
\end{itemize}

Hence:
\[
c_0=1,\qquad c_1=2g,\qquad c_2=1.
\]
So
\[
\chi(\Sigma_g)=1-2g+1=2-2g.
\]

\subsubsection*{Examples}

For \(g=1\):
\[
a_1b_1a_1^{-1}b_1^{-1},
\]
which is the torus, and
\[
\chi(\Sigma_1)=2-2=0.
\]

For \(g=2\):
\[
a_1b_1a_1^{-1}b_1^{-1}a_2b_2a_2^{-1}b_2^{-1},
\]
so
\[
\chi(\Sigma_2)=2-4=-2.
\]

Each additional handle adds two new \(1\)-cells and keeps \(c_0\) and \(c_2\) the same, so the Euler characteristic drops by \(2\).

\subsection{Nonorientable surface of genus \(n\) from a \(2n\)-gon}

A standard polygonal model for the nonorientable surface \(S_n\) is:
\[
a_1a_1a_2a_2\cdots a_n a_n.
\]

This is a \(2n\)-gon, with each pair of matching edges glued in the same direction.

\subsubsection*{Cell count}

In the quotient:
\begin{itemize}
	\item all vertices identify to one vertex,
	\item each pair of edges gives one \(1\)-cell, so there are \(n\) such cells,
	\item the interior gives one \(2\)-cell.
\end{itemize}

Thus:
\[
c_0=1,\qquad c_1=n,\qquad c_2=1.
\]
Hence
\[
\chi(S_n)=1-n+1=2-n.
\]

\subsubsection*{Examples}

For \(n=1\):
\[
aa,
\]
which gives \(\RP^2\), and
\[
\chi(S_1)=2-1=1.
\]

For \(n=2\):
\[
aabb,
\]
which gives the Klein bottle, and
\[
\chi(S_2)=2-2=0.
\]

For \(n=3\):
\[
aabbcc,
\]
which gives \(S_3\), and
\[
\chi(S_3)=2-3=-1.
\]

Each additional crosscap adds one more \(1\)-cell and lowers the Euler characteristic by \(1\).

\subsection{Reading Euler characteristic directly from the polygon model}

For the standard one-polygon models above, the computation is extremely simple:
\begin{itemize}
	\item there is one \(2\)-cell: the polygon interior,
	\item all vertices collapse to one \(0\)-cell,
	\item the number of \(1\)-cells is the number of edge-pairs.
\end{itemize}

So:
\[
\chi = 1 - (\text{number of edge-pairs}) + 1.
\]

Therefore:
\begin{itemize}
	\item in the orientable genus-\(g\) case there are \(2g\) edge-pairs, so
	\[
	\chi=2-2g,
	\]
	\item in the nonorientable genus-\(n\) case there are \(n\) edge-pairs, so
	\[
	\chi=2-n.
	\]
\end{itemize}

This is one of the cleanest ways to remember the formulas.

\subsection{A useful comparison table}

\[
\begin{array}{c|c|c|c|c}
	\text{Surface} & \text{Polygon word} & c_0 & c_1 & \chi \\ \hline
	S^2 & \text{(triangulated sphere)} & - & - & 2 \\
	T^2 & aba^{-1}b^{-1} & 1 & 2 & 0 \\
	\RP^2 & aa & 1 & 1 & 1 \\
	K & abab & 1 & 2 & 0 \\
	\Sigma_g & a_1b_1a_1^{-1}b_1^{-1}\cdots a_gb_ga_g^{-1}b_g^{-1} & 1 & 2g & 2-2g \\
	S_n & a_1a_1a_2a_2\cdots a_na_n & 1 & n & 2-n
\end{array}
\]

\subsection{Why polygonal models are so important}

Polygonal models are useful because they connect:
\begin{itemize}
	\item geometry: draw a polygon and glue edges,
	\item combinatorics: count vertices, edges, and faces,
	\item topology: identify the resulting surface,
	\item algebra: later read off fundamental groups and homology presentations.
\end{itemize}

They are one of the best bridges between concrete pictures and abstract invariants.

\subsection{Summary}

The polygonal model of a surface gives one of the fastest routes to the Euler characteristic.

For the torus:
\[
aba^{-1}b^{-1}\quad \Longrightarrow \quad c_0=1,\ c_1=2,\ c_2=1,\ \chi=0.
\]

For the projective plane:
\[
aa\quad \Longrightarrow \quad c_0=1,\ c_1=1,\ c_2=1,\ \chi=1.
\]

For the Klein bottle:
\[
abab\quad \Longrightarrow \quad c_0=1,\ c_1=2,\ c_2=1,\ \chi=0.
\]

In general:
\[
\chi(\Sigma_g)=2-2g,\qquad \chi(S_n)=2-n.
\]

So polygonal edge-identification models are not just pictures: they are concrete computational tools for topology.

\section{From polygon words to presentations of the fundamental group}

A polygonal model of a surface does not only allow us to compute the Euler characteristic. It also gives a direct presentation of the fundamental group. This is one of the most important bridges between geometry and algebra in topology.

\subsection{General principle}

Suppose a closed surface is obtained from a polygon whose boundary is labeled by a cyclic word
\[
w
\]
in symbols \(a,b,c,\dots\), with appropriate inverses when orientations are reversed.

After the edge identifications:
\begin{itemize}
	\item all vertices often collapse to a single \(0\)-cell,
	\item each pair of identified edges gives a \(1\)-cell,
	\item the interior of the polygon becomes one \(2\)-cell.
\end{itemize}

Thus the quotient has a CW decomposition with:
\begin{itemize}
	\item one \(0\)-cell,
	\item several \(1\)-cells,
	\item one \(2\)-cell attached along the boundary word \(w\).
\end{itemize}

Therefore:
\begin{itemize}
	\item the \(1\)-skeleton is a wedge of circles,
	\item its fundamental group is a free group on the edge labels,
	\item attaching the \(2\)-cell imposes one relation: the boundary word becomes null-homotopic.
\end{itemize}

So the resulting presentation is
\[
\pi_1(X)\cong \langle \text{edge labels}\mid w=1\rangle.
\]

\subsection{Why this works}

Before attaching the interior of the polygon, the quotient space is just a graph obtained from the identified boundary edges. Since all edges meet at the same vertex, this graph is a wedge of circles.

If there are \(r\) edge-pairs, then the \(1\)-skeleton is homotopy equivalent to
\[
S^1\vee S^1\vee \cdots \vee S^1
\]
(\(r\) copies), so its fundamental group is the free group
\[
F_r.
\]

Now attach the polygon interior as a \(2\)-cell. The boundary of that \(2\)-cell travels around the \(1\)-skeleton according to the boundary word \(w\). Since the boundary of a \(2\)-cell is null-homotopic in the attached space, the word \(w\) becomes the relation.

This is an application of Seifert--van Kampen's theorem.

\subsection{Warm-up: the wedge of circles}

Let
\[
X=S^1\vee S^1
\]
be the wedge of two circles with loops labeled \(a\) and \(b\). Then
\[
\pi_1(X)\cong F_2=\langle a,b\rangle.
\]

If one now attaches a \(2\)-cell along the loop
\[
aba^{-1}b^{-1},
\]
then the resulting space has fundamental group
\[
\langle a,b\mid aba^{-1}b^{-1}=1\rangle.
\]

This is exactly the torus case.

\subsection{Example 1: the torus}

The torus \(T^2\) is obtained from a square with boundary word
\[
aba^{-1}b^{-1}.
\]

Hence
\[
\pi_1(T^2)\cong \langle a,b\mid aba^{-1}b^{-1}=1\rangle.
\]

Using commutator notation
\[
[a,b]=aba^{-1}b^{-1},
\]
we may also write
\[
\pi_1(T^2)\cong \langle a,b\mid [a,b]=1\rangle.
\]

This means that \(a\) and \(b\) commute, so
\[
\pi_1(T^2)\cong \Z\oplus \Z.
\]

\subsubsection*{Geometric meaning}

The loop \(a\) goes once around one basic direction of the torus, and \(b\) goes once around the other. The relation says that going around \(a\) and then \(b\) is homotopic to going around \(b\) and then \(a\).

So the torus has abelian fundamental group.

\subsection{Example 2: the projective plane}

A standard polygonal model for \(\RP^2\) has boundary word
\[
aa.
\]

Thus
\[
\pi_1(\RP^2)\cong \langle a\mid a^2=1\rangle.
\]

Therefore
\[
\pi_1(\RP^2)\cong \Z/2.
\]

\subsubsection*{Geometric meaning}

There is one basic nontrivial loop \(a\), but traversing it twice becomes null-homotopic. That is exactly what the relation
\[
a^2=1
\]
expresses.

\subsection{Example 3: the Klein bottle}

A standard polygonal model for the Klein bottle has boundary word
\[
abab.
\]

Hence
\[
\pi_1(K)\cong \langle a,b\mid abab=1\rangle.
\]

This can be rewritten in a more familiar form. From
\[
abab=1
\]
we get
\[
ab=b^{-1}a^{-1},
\]
and hence
\[
aba^{-1}=b^{-1}.
\]

So an equivalent presentation is
\[
\pi_1(K)\cong \langle a,b\mid aba^{-1}=b^{-1}\rangle.
\]

\subsubsection*{Geometric meaning}

The generator \(a\) acts on \(b\) by reversing it. Thus the group is nonabelian:
\[
ab\neq ba.
\]

This algebraic noncommutativity reflects the nonorientable geometry of the Klein bottle.

\subsection{Example 4: orientable surface of genus \(g\)}

The standard polygonal model for the orientable surface \(\Sigma_g\) has boundary word
\[
a_1b_1a_1^{-1}b_1^{-1}\cdots a_gb_ga_g^{-1}b_g^{-1}.
\]

Therefore
\[
\pi_1(\Sigma_g)\cong
\left\langle
a_1,b_1,\dots,a_g,b_g
\;\middle|\;
a_1b_1a_1^{-1}b_1^{-1}\cdots a_gb_ga_g^{-1}b_g^{-1}=1
\right\rangle.
\]

Using commutators
\[
[x,y]=xyx^{-1}y^{-1},
\]
this becomes
\[
\pi_1(\Sigma_g)\cong
\left\langle
a_1,b_1,\dots,a_g,b_g
\;\middle|\;
[a_1,b_1]\cdots [a_g,b_g]=1
\right\rangle.
\]

This is the standard presentation of the orientable surface group.

\subsubsection*{Special cases}

For \(g=1\),
\[
\pi_1(\Sigma_1)=\pi_1(T^2)\cong \langle a,b\mid [a,b]=1\rangle \cong \Z^2.
\]

For \(g=2\),
\[
\pi_1(\Sigma_2)\cong
\langle a_1,b_1,a_2,b_2\mid [a_1,b_1][a_2,b_2]=1\rangle.
\]

For \(g\ge 2\), this group is nonabelian.

\subsection{Example 5: nonorientable surface of genus \(n\)}

The standard polygonal model for the nonorientable surface \(S_n\) has boundary word
\[
a_1a_1a_2a_2\cdots a_n a_n.
\]

Hence
\[
\pi_1(S_n)\cong
\left\langle
a_1,\dots,a_n
\;\middle|\;
a_1^2a_2^2\cdots a_n^2=1
\right\rangle.
\]

This is the standard presentation of the nonorientable surface group.

\subsubsection*{Special cases}

For \(n=1\),
\[
\pi_1(S_1)=\pi_1(\RP^2)\cong \langle a\mid a^2=1\rangle \cong \Z/2.
\]

For \(n=2\),
\[
\pi_1(S_2)\cong \langle a,b\mid a^2b^2=1\rangle.
\]

Since \(S_2\) is the Klein bottle, this presentation is equivalent to the usual Klein bottle presentation.

\subsection{Why the boundary word gives the relation}

This is the central geometric idea.

Before attaching the polygon interior, the identified boundary gives a loop in the \(1\)-skeleton. That loop spells out exactly the boundary word \(w\).

When the polygon interior is attached, that boundary loop becomes the boundary of a disk. Any loop bounding a disk is null-homotopic. Therefore the corresponding word must equal the identity in the fundamental group.

So the boundary word becomes the defining relation.

\subsection{Relation with van Kampen's theorem}

This method is a particularly concrete instance of Seifert--van Kampen's theorem.

The general pattern is:
\begin{itemize}
	\item the \(1\)-skeleton contributes a free group,
	\item the \(2\)-cell kills one loop,
	\item therefore the fundamental group is the free group modulo the normal closure of the boundary word.
\end{itemize}

This is exactly why polygonal models are so effective.

\subsection{Passing from \(\pi_1\) to \(H_1\) by abelianization}

A very important fact is:
\[
H_1(X;\Z)\cong \pi_1(X)^{ab},
\]
the abelianization of the fundamental group.

That means we obtain \(H_1\) by forcing all generators of \(\pi_1\) to commute.

\subsubsection*{Torus}

From
\[
\pi_1(T^2)\cong \langle a,b\mid aba^{-1}b^{-1}=1\rangle,
\]
the relation becomes trivial after abelianization, since commutators vanish automatically. Thus
\[
H_1(T^2;\Z)\cong \Z^2.
\]

\subsubsection*{Projective plane}

From
\[
\pi_1(\RP^2)\cong \langle a\mid a^2=1\rangle,
\]
we obtain
\[
H_1(\RP^2;\Z)\cong \Z/2.
\]

\subsubsection*{Klein bottle}

From
\[
\pi_1(K)\cong \langle a,b\mid aba^{-1}=b^{-1}\rangle,
\]
abelianization gives
\[
a+b-a=-b,
\]
so
\[
2b=0.
\]
Hence
\[
H_1(K;\Z)\cong \Z\oplus \Z/2.
\]

\subsubsection*{Orientable surface of genus \(g\)}

From
\[
\pi_1(\Sigma_g)\cong
\left\langle
a_1,b_1,\dots,a_g,b_g
\;\middle|\;
[a_1,b_1]\cdots [a_g,b_g]=1
\right\rangle,
\]
all commutators disappear after abelianization, so
\[
H_1(\Sigma_g;\Z)\cong \Z^{2g}.
\]

\subsubsection*{Nonorientable surface of genus \(n\)}

From
\[
\pi_1(S_n)\cong
\left\langle
a_1,\dots,a_n
\;\middle|\;
a_1^2\cdots a_n^2=1
\right\rangle,
\]
abelianization gives the additive relation
\[
2a_1+\cdots+2a_n=0.
\]
Thus
\[
H_1(S_n;\Z)\cong \Z^{n-1}\oplus \Z/2.
\]

So the polygon presentation gives both the fundamental group and the first homology.

\subsection{A compact summary of the main presentations}

\[
\pi_1(T^2)\cong \langle a,b\mid aba^{-1}b^{-1}=1\rangle \cong \Z^2,
\]

\[
\pi_1(\RP^2)\cong \langle a\mid a^2=1\rangle \cong \Z/2,
\]

\[
\pi_1(K)\cong \langle a,b\mid aba^{-1}=b^{-1}\rangle,
\]

\[
\pi_1(\Sigma_g)\cong
\left\langle
a_1,b_1,\dots,a_g,b_g
\;\middle|\;
[a_1,b_1]\cdots [a_g,b_g]=1
\right\rangle,
\]

\[
\pi_1(S_n)\cong
\left\langle
a_1,\dots,a_n
\;\middle|\;
a_1^2a_2^2\cdots a_n^2=1
\right\rangle.
\]

\subsection{Summary}

A polygonal model of a surface immediately gives:
\begin{itemize}
	\item the Euler characteristic by counting cells,
	\item the fundamental group by taking generators from the edge-pairs and one relation from the boundary word,
	\item the first homology by abelianizing the resulting group presentation.
\end{itemize}

Thus one polygon picture contains a remarkable amount of topology:
\begin{itemize}
	\item geometry of the gluing,
	\item combinatorics of the cell structure,
	\item algebra of the fundamental group,
	\item homology of the surface.
\end{itemize}

\section{From presentations to explicit generators of \(H_1\), loops on the surface, and the geometric meaning of commutators and squares}

The polygonal presentation of a surface gives a presentation of its fundamental group. The next step is to understand what this says geometrically about loops on the surface, and how one passes from those loops to explicit generators of the first homology group
\[
H_1(X;\Z).
\]

This section explains:
\begin{itemize}
	\item how the edge loops in a polygon model give concrete loops on the surface,
	\item how \(H_1\) is obtained from \(\pi_1\) by abelianization,
	\item what the generators of \(H_1\) look like geometrically,
	\item why commutators disappear in \(H_1\),
	\item why squares in nonorientable surfaces produce the \(\Z/2\) torsion.
\end{itemize}

\subsection{From edge labels to actual loops on the surface}

Start with a polygonal model of a surface in which all vertices are identified to a single point \(x_0\). After the edge identifications, each edge-pair determines a loop based at \(x_0\).

So if the polygon has edge labels
\[
a,b,c,\dots,
\]
then on the quotient surface these become based loops
\[
\alpha,\beta,\gamma,\dots
\]
starting and ending at the same point \(x_0\).

These loops are the geometric representatives of the generators in the presentation of the fundamental group.

\subsubsection*{Example: the torus}

For the square model of the torus with word
\[
aba^{-1}b^{-1},
\]
the two edge-pairs become two loops on the torus:
\begin{itemize}
	\item one loop going once around the torus in the horizontal direction,
	\item one loop going once around the torus in the vertical direction.
\end{itemize}

These are the standard two generators of
\[
\pi_1(T^2)\cong \Z^2.
\]

\subsubsection*{Example: the projective plane}

For the model
\[
aa,
\]
there is one loop \(a\). Geometrically, this is the basic nontrivial loop in \(\RP^2\). It is not null-homotopic, but going around it twice becomes null-homotopic.

\subsubsection*{Example: the Klein bottle}

For the model
\[
abab,
\]
there are two basic loops \(a\) and \(b\), but their interaction is different from the torus case: one of them reverses the other under conjugation.

\subsection{From \(\pi_1\) to \(H_1\): abelianization}

The first homology group is obtained from the fundamental group by abelianization:
\[
H_1(X;\Z)\cong \pi_1(X)^{ab}.
\]

This means:
\begin{itemize}
	\item take the group \(\pi_1(X)\),
	\item force all generators to commute,
	\item interpret the resulting group additively.
\end{itemize}

So if \(\pi_1(X)\) is written multiplicatively in generators \(a,b,\dots\), then in \(H_1(X;\Z)\) we usually write the corresponding homology classes additively:
\[
[a],[b],\dots
\]
or simply
\[
a,b,\dots
\]
when no confusion can occur.

In homology:
\[
ab=ba,
\]
so multiplicative group words become additive relations.

For example:
\[
aba^{-1}b^{-1}=1
\]
becomes
\[
a+b-a-b=0.
\]

\subsection{Why \(H_1\) forgets order}

The fundamental group remembers the order in which loops are traversed. Thus in general
\[
ab\neq ba
\]
in \(\pi_1(X)\).

But \(H_1\) only records the \(1\)-dimensional cycle carried by the loop. At that level, order no longer matters. Traversing loop \(a\) then \(b\) gives the same \(1\)-cycle as traversing \(b\) then \(a\).

That is why \(H_1\) is abelian.

Geometrically, \(H_1\) forgets the fine noncommutative structure of how loops wind; it remembers only their total \(1\)-dimensional homology classes.

\subsection{The geometric meaning of a commutator}

Recall that the commutator of two loops is
\[
[a,b]=aba^{-1}b^{-1}.
\]

Geometrically, this means:
\begin{itemize}
	\item go along \(a\),
	\item then along \(b\),
	\item then back along \(a\),
	\item then back along \(b\).
\end{itemize}

On the torus, this loop goes once around the boundary of the fundamental square. But the interior of that square becomes a \(2\)-cell of the torus. Therefore the commutator loop bounds that \(2\)-cell, so it is trivial in \(\pi_1(T^2)\).

In homology, every commutator disappears automatically:
\[
[a,b]\mapsto 0 \in H_1(X;\Z).
\]

So commutators represent the part of the fundamental group that is invisible to \(H_1\).

\subsubsection*{Torus picture}

For the torus:
\[
\pi_1(T^2)\cong \langle a,b\mid aba^{-1}b^{-1}=1\rangle.
\]
The relation says that the commutator bounds the \(2\)-cell. Thus in homology, the classes \(a\) and \(b\) survive freely, and there is no further relation:
\[
H_1(T^2;\Z)\cong \Z a\oplus \Z b\cong \Z^2.
\]

\subsection{The geometric meaning of a square}

In nonorientable surfaces, the basic relation often involves squares:
\[
a_1^2a_2^2\cdots a_n^2=1.
\]

A square \(a^2\) means traversing the same loop twice in succession.

This has a geometric meaning in nonorientable surfaces: a loop passing once through a crosscap reverses local orientation, and in the cell decomposition the boundary of the \(2\)-cell runs twice along that edge. That is why the relation contains a square.

So:
\begin{itemize}
	\item commutators correspond to orientable handle structure,
	\item squares correspond to crosscaps and nonorientable structure.
\end{itemize}

\subsubsection*{Projective plane}

In \(\RP^2\), the relation is
\[
a^2=1.
\]
After abelianization this becomes
\[
2a=0.
\]
Thus
\[
H_1(\RP^2;\Z)\cong \Z/2.
\]

Geometrically, the basic loop \(a\) is nontrivial, but twice around it bounds the \(2\)-cell. So its homology class has order \(2\).

\subsection{Explicit generators of \(H_1(T^2;\Z)\)}

The torus has presentation
\[
\pi_1(T^2)\cong \langle a,b\mid aba^{-1}b^{-1}=1\rangle.
\]

After abelianization, the commutator relation disappears, so
\[
H_1(T^2;\Z)\cong \Z^2.
\]

The explicit generators are:
\begin{itemize}
	\item the homology class \([a]\) of the loop going once around the torus meridionally,
	\item the homology class \([b]\) of the loop going once around the torus longitudinally.
\end{itemize}

Every homology class in \(H_1(T^2;\Z)\) is of the form
\[
m[a]+n[b],\qquad m,n\in \Z.
\]

Geometrically, such a class is represented by a loop winding \(m\) times in one basic direction and \(n\) times in the other.

\subsection{Explicit generators of \(H_1(\Sigma_g;\Z)\)}

For the orientable surface of genus \(g\),
\[
\pi_1(\Sigma_g)\cong
\left\langle
a_1,b_1,\dots,a_g,b_g
\;\middle|\;
[a_1,b_1]\cdots[a_g,b_g]=1
\right\rangle.
\]

After abelianization, all commutators vanish automatically, so no nontrivial relation remains among the generators. Hence
\[
H_1(\Sigma_g;\Z)\cong \Z^{2g}.
\]

An explicit basis is:
\[
[a_1],[b_1],\dots,[a_g],[b_g].
\]

Geometrically:
\begin{itemize}
	\item each handle contributes two independent loops,
	\item one goes around the handle in one direction,
	\item the other goes around it in a transverse direction.
\end{itemize}

Thus each handle contributes two free homology generators.

\subsubsection*{Example: genus \(2\)}

For \(\Sigma_2\),
\[
\pi_1(\Sigma_2)\cong \langle a_1,b_1,a_2,b_2\mid [a_1,b_1][a_2,b_2]=1\rangle,
\]
so
\[
H_1(\Sigma_2;\Z)\cong \Z^4
\]
with basis
\[
[a_1],[b_1],[a_2],[b_2].
\]

\subsection{Explicit generators of \(H_1(\RP^2;\Z)\)}

For the projective plane,
\[
\pi_1(\RP^2)\cong \langle a\mid a^2=1\rangle.
\]

Abelianization changes this into
\[
2[a]=0.
\]

So
\[
H_1(\RP^2;\Z)\cong \Z/2,
\]
generated by the class \([a]\).

This generator is torsion: it is nonzero, but
\[
2[a]=0.
\]

Geometrically, this says:
\begin{itemize}
	\item one pass around the basic nontrivial loop gives a nontrivial homology class,
	\item two passes give a boundary.
\end{itemize}

\subsection{Explicit generators of \(H_1(K;\Z)\)}

For the Klein bottle,
\[
\pi_1(K)\cong \langle a,b\mid aba^{-1}=b^{-1}\rangle.
\]

Abelianization turns this into
\[
a+b-a=-b,
\]
hence
\[
2b=0.
\]

Therefore
\[
H_1(K;\Z)\cong \Z\oplus \Z/2.
\]

An explicit choice of generators is:
\begin{itemize}
	\item \([a]\), which generates the free \(\Z\)-part,
	\item \([b]\), which generates the torsion \(\Z/2\)-part.
\end{itemize}

So every class in \(H_1(K;\Z)\) has the form
\[
m[a]+\varepsilon [b],\qquad m\in \Z,\ \varepsilon\in\{0,1\}.
\]

\subsubsection*{Geometric meaning}

One basic direction on the Klein bottle behaves like an ordinary circle and gives an infinite-order homology class. The other direction is twisted by the nonorientable gluing, and its homology class has order \(2\).

This is a very concrete place where torsion appears geometrically.

\subsection{Explicit generators of \(H_1(S_n;\Z)\)}

For the nonorientable surface \(S_n\),
\[
\pi_1(S_n)\cong
\left\langle
a_1,\dots,a_n
\;\middle|\;
a_1^2a_2^2\cdots a_n^2=1
\right\rangle.
\]

After abelianization, the relation becomes
\[
2[a_1]+\cdots+2[a_n]=0.
\]

So
\[
H_1(S_n;\Z)\cong \Z^{n-1}\oplus \Z/2.
\]

A convenient way to see generators is this:
\begin{itemize}
	\item differences such as
	\[
	[a_1]-[a_n],\ [a_2]-[a_n],\ \dots,\ [a_{n-1}]-[a_n]
	\]
	generate the free part \(\Z^{n-1}\),
	\item the sum
	\[
	[a_1]+\cdots+[a_n]
	\]
	gives the torsion class of order \(2\).
\end{itemize}

\subsubsection*{Why the torsion has order \(2\)}

From
\[
2[a_1]+\cdots+2[a_n]=0
\]
we get
\[
2([a_1]+\cdots+[a_n])=0.
\]
So the total sum is a \(2\)-torsion class.

This is the algebraic shadow of the nonorientable gluing.

\subsection{Concrete comparison: handle versus crosscap}

There is a very important geometric contrast.

\subsubsection*{Handle}

A handle contributes a commutator relation:
\[
[a,b]=aba^{-1}b^{-1}.
\]
In homology this disappears, so both \(a\) and \(b\) survive as free generators. Thus one handle contributes
\[
\Z^2
\]
to \(H_1\).

\subsubsection*{Crosscap}

A crosscap contributes a square relation:
\[
a^2.
\]
In homology this becomes
\[
2a=0
\]
or, in the many-crosscap case, a sum of doubled generators. Thus crosscaps are the source of \(2\)-torsion.

So:
\begin{itemize}
	\item handles create free generators,
	\item crosscaps create \(2\)-torsion.
\end{itemize}

\subsection{Loops, orientations, and homology classes}

If \(\gamma\) is a loop, then reversing its direction changes the homology class to its negative:
\[
[\gamma^{-1}]=-[\gamma].
\]

So in additive notation:
\[
a^{-1} \leadsto -a.
\]

That is why a word such as
\[
aba^{-1}b^{-1}
\]
becomes
\[
a+b-a-b=0
\]
in homology.

This passage from multiplicative notation in \(\pi_1\) to additive notation in \(H_1\) is one of the most important computational tools in low-dimensional topology.

\subsection{A chain-level viewpoint}

There is also a homological way to see the same story.

In a CW decomposition:
\begin{itemize}
	\item the \(1\)-cells generate \(C_1(X;\Z)\),
	\item the attaching map of the \(2\)-cell determines the boundary homomorphism
	\[
	\partial_2:C_2(X;\Z)\to C_1(X;\Z).
	\]
\end{itemize}

Then
\[
H_1(X;\Z)=\ker(\partial_1)/\operatorname{im}(\partial_2).
\]

Since there is only one vertex, usually \(\partial_1=0\), so
\[
H_1(X;\Z)\cong C_1(X;\Z)/\operatorname{im}(\partial_2).
\]

Thus:
\begin{itemize}
	\item the \(1\)-cells give generators,
	\item the boundary of the \(2\)-cell gives the relation.
\end{itemize}

This exactly matches the abelianization of the fundamental group presentation.

\subsection{A compact table of generators and relations in \(H_1\)}

\[
\begin{array}{c|c|c}
	X & \text{Generators in }H_1 & \text{Relation in }H_1 \\ \hline
	T^2 & [a],[b] & \text{none} \\
	\RP^2 & [a] & 2[a]=0 \\
	K & [a],[b] & 2[b]=0 \\
	\Sigma_g & [a_1],[b_1],\dots,[a_g],[b_g] & \text{none} \\
	S_n & [a_1],\dots,[a_n] & 2[a_1]+\cdots+2[a_n]=0
\end{array}
\]

So:
\[
H_1(T^2;\Z)\cong \Z^2,
\qquad
H_1(\RP^2;\Z)\cong \Z/2,
\qquad
H_1(K;\Z)\cong \Z\oplus \Z/2,
\]
\[
H_1(\Sigma_g;\Z)\cong \Z^{2g},
\qquad
H_1(S_n;\Z)\cong \Z^{n-1}\oplus \Z/2.
\]

\subsection{Summary}

The polygon presentation of a surface gives concrete edge loops, and these loops become the basic generators of both \(\pi_1\) and \(H_1\).

The key principles are:
\begin{itemize}
	\item edge-pairs give loops on the surface,
	\item \(H_1\) is obtained from \(\pi_1\) by abelianization,
	\item commutators vanish in \(H_1\),
	\item squares produce \(2\)-torsion,
	\item handles contribute free generators,
	\item crosscaps contribute torsion phenomena.
\end{itemize}

Thus one passes from the geometric picture of loops on the surface to explicit algebraic generators of homology in a direct and computable way.

\section{What is a \(2\)-cycle, and why does \(S^2\) have a nontrivial \(2\)-dimensional homology class?}

A very common source of confusion is the phrase ``the sphere has one \(2\)-cycle.''  
Strictly speaking, this is not the best wording. The correct statement is:
\[
H_2(S^2;\Z)\cong \Z.
\]
This means that \(S^2\) has one \emph{independent generator} in dimension \(2\), not that there is literally only one \(2\)-cycle.

\subsection{Chains, cycles, and boundaries}

Let \(K\) be a simplicial complex. A simplicial \(2\)-chain is a finite formal integer linear combination of oriented \(2\)-simplices:
\[
c=n_1\sigma_1+n_2\sigma_2+\cdots+n_r\sigma_r,
\qquad n_i\in \Z.
\]

Here each \(\sigma_i\) is an oriented triangle.

The boundary map
\[
\partial_2:C_2(K;\Z)\to C_1(K;\Z)
\]
sends each oriented triangle to the oriented sum of its three edges. For an oriented simplex
\[
[v_0,v_1,v_2],
\]
the formula is
\[
\partial [v_0,v_1,v_2]
=
[v_1,v_2]-[v_0,v_2]+[v_0,v_1].
\]

\begin{definition}
	A \(2\)-chain \(c\) is called a \emph{\(2\)-cycle} if
	\[
	\partial_2(c)=0.
	\]
	That is, the boundaries of all the triangles cancel.
\end{definition}

So geometrically, a \(2\)-cycle is a formal sum of oriented triangles which closes up with no leftover boundary edges.

\begin{definition}
	A \(2\)-cycle \(c\) is called a \emph{\(2\)-boundary} if there exists a \(3\)-chain \(d\) such that
	\[
	c=\partial_3(d).
	\]
\end{definition}

The second homology group is then
\[
H_2(K;\Z)=Z_2(K;\Z)/B_2(K;\Z),
\]
that is, \(2\)-cycles modulo \(2\)-boundaries.

\subsection{The basic geometric picture}

It is useful to think in dimensions:

\begin{itemize}
	\item a \(0\)-cycle is a formal sum of points,
	\item a \(1\)-cycle is a closed loop,
	\item a \(2\)-cycle is a closed shell or closed surface made from triangles.
\end{itemize}

Thus for \(S^2\), the whole sphere itself is the basic geometric model of a \(2\)-cycle.

\subsection{The sphere as the boundary of a tetrahedron}

A very concrete triangulation of \(S^2\) is the boundary of a tetrahedron. It consists of \(4\) triangular faces:
\[
\sigma_1,\sigma_2,\sigma_3,\sigma_4.
\]

If these faces are oriented consistently, for example all outward, then the sum
\[
c=\sigma_1+\sigma_2+\sigma_3+\sigma_4
\]
is a \(2\)-chain.

Now compute its boundary:
\[
\partial c
=
\partial\sigma_1+\partial\sigma_2+\partial\sigma_3+\partial\sigma_4.
\]

Each edge of the tetrahedron belongs to exactly two faces, and those two faces induce opposite orientations on that common edge. Therefore every edge appears twice with opposite signs, so all edge terms cancel.

Hence
\[
\partial c=0.
\]

Therefore \(c\) is a \(2\)-cycle.

\subsection{Why this \(2\)-cycle is nontrivial}

The crucial point is that this \(2\)-cycle is not a boundary inside \(S^2\).

If the sphere were sitting as the boundary of a \(3\)-dimensional simplex or a \(3\)-chain \emph{inside the space}, then it would represent zero in homology. But the space here is only \(S^2\) itself. It contains triangles, edges, and vertices in the triangulation, but no \(3\)-simplices.

So there is no \(3\)-chain \(d\) in \(S^2\) such that
\[
\partial_3(d)=c.
\]

Thus \(c\) is a \(2\)-cycle that is not a \(2\)-boundary, and hence it determines a nonzero class in
\[
H_2(S^2;\Z).
\]

\subsection{Why the homology group is \(\Z\)}

The class of the whole oriented sphere is called the \emph{fundamental class} of \(S^2\), often written
\[
[S^2].
\]

Every integer multiple
\[
n[S^2],\qquad n\in \Z,
\]
is again a homology class in dimension \(2\). Reversing the orientation changes the sign:
\[
-[S^2].
\]

So the second homology group has the form
\[
H_2(S^2;\Z)\cong \Z.
\]

This means:
\begin{itemize}
	\item \(0\) corresponds to the trivial class,
	\item \(1\) corresponds to the positively oriented sphere,
	\item \(-1\) corresponds to the oppositely oriented sphere,
	\item \(2,3,\dots\) correspond to multiple copies.
\end{itemize}

\subsection{Important correction of language}

It is better to say:
\[
\text{``\(S^2\) has one generator in \(H_2\)''}
\]
rather than
\[
\text{``\(S^2\) has one \(2\)-cycle.''}
\]

Indeed, \(S^2\) has many \(2\)-cycles:
\[
c,\quad 2c,\quad -c,\quad 3c,\dots
\]
and many different representatives of the same homology class. What homology says is that all \(2\)-dimensional classes are integer multiples of one basic class.

\subsection{Comparison with \(S^1\)}

This also explains the difference between the circle and the sphere.

For the circle \(S^1\),
\[
H_1(S^1;\Z)\cong \Z,
\qquad
H_2(S^1;\Z)=0.
\]

The circle has a nontrivial \(1\)-cycle, namely the loop itself, but it has no \(2\)-dimensional shell.

For the sphere \(S^2\),
\[
H_2(S^2;\Z)\cong \Z.
\]

The sphere has a nontrivial \(2\)-cycle, namely the whole oriented surface.

\subsection{A useful geometric summary}

A \(2\)-cycle is an oriented sum of triangles with no boundary:
\[
\partial_2(c)=0.
\]

For \(S^2\), the whole sphere is such a \(2\)-cycle. In a triangulation, it is the sum of all faces with consistent orientation. Every edge cancels, so the boundary is zero.

This \(2\)-cycle is not a boundary inside \(S^2\), because \(S^2\) contains no \(3\)-dimensional chain whose boundary is the whole sphere.

Therefore it determines a nonzero homology class, and in fact
\[
H_2(S^2;\Z)\cong \Z.
\]

\subsection{Compact summary}

\begin{itemize}
	\item A \(2\)-chain is a formal sum of oriented triangles.
	\item A \(2\)-cycle is a \(2\)-chain with zero boundary.
	\item On \(S^2\), the sum of all triangles in a triangulation gives a \(2\)-cycle.
	\item This cycle is not a boundary in \(S^2\).
	\item Therefore it gives a nontrivial class in \(H_2(S^2;\Z)\).
	\item All classes are integer multiples of one generator, so
	\[
	H_2(S^2;\Z)\cong \Z.
	\]
\end{itemize}

\section{Picture-style explanation of cycles vs boundaries in dimensions \(0\), \(1\), and \(2\)}

The central idea of homology is very simple:

\begin{itemize}
	\item a \emph{cycle} is something that has no boundary left over,
	\item a \emph{boundary} is something that actually arises as the boundary of a higher-dimensional object.
\end{itemize}

Thus in each dimension one asks two questions:
\begin{enumerate}
	\item does it close up?
	\item does it bound something one dimension higher?
\end{enumerate}

Homology measures cycles modulo boundaries:
\[
H_n(X)=Z_n(X)/B_n(X).
\]

So the philosophy is:
\begin{itemize}
	\item first find all \(n\)-cycles,
	\item then declare every \(n\)-boundary to be trivial,
	\item what remains are the genuine \(n\)-dimensional holes of the space.
\end{itemize}

\subsection{Dimension \(0\): points}

A \(0\)-chain is a formal integer sum of points:
\[
c = n_1p_1+n_2p_2+\cdots+n_rp_r.
\]

Since points have no lower-dimensional boundary, every \(0\)-chain automatically has zero boundary. Thus:
\[
Z_0(X)=C_0(X).
\]

So every \(0\)-chain is a \(0\)-cycle.

\subsubsection*{What is a \(0\)-boundary?}

A \(0\)-boundary comes from the boundary of a \(1\)-chain, that is, from an edge or sum of edges.

If \([p,q]\) is an oriented edge from \(p\) to \(q\), then
\[
\partial[p,q]=q-p.
\]

So the difference of two points is a boundary whenever they are joined by an edge, or more generally by a path in the same connected component.

\subsubsection*{Picture intuition}

In dimension \(0\):
\begin{itemize}
	\item cycles are just collections of points,
	\item boundaries are differences of endpoints of paths,
	\item homology detects connected components.
\end{itemize}

If two points lie in the same connected component, then their difference is a boundary, so in homology they represent the same class.

\subsubsection*{Example 1: a connected space}

If \(X\) is connected, then all points are homologous, so
\[
H_0(X)\cong \Z.
\]

There is one generator corresponding to the one connected component.

\subsubsection*{Example 2: two disconnected points}

If \(X=\{p,q\}\) has two connected components, then there is no edge or path joining them, so \(q-p\) is not a boundary. Hence
\[
H_0(X)\cong \Z^2.
\]

One copy of \(\Z\) corresponds to each connected component.

\subsection{Dimension \(1\): loops}

A \(1\)-chain is a formal integer sum of oriented edges:
\[
c=n_1e_1+n_2e_2+\cdots+n_re_r.
\]

A \(1\)-cycle is a \(1\)-chain whose boundary is zero:
\[
\partial_1(c)=0.
\]

This means that the endpoints cancel, so there is no leftover start or endpoint.

\subsubsection*{Picture intuition}

A \(1\)-cycle is essentially a closed loop, or a union of closed loops.

Examples include:
\begin{itemize}
	\item a circle,
	\item the boundary of a triangle,
	\item the boundary of a square,
	\item several disjoint loops.
\end{itemize}

\subsubsection*{What is a \(1\)-boundary?}

A \(1\)-boundary is the boundary of a \(2\)-chain, that is, the boundary of some collection of oriented triangles or filled-in surface pieces.

For an oriented triangle \([v_0,v_1,v_2]\),
\[
\partial[v_0,v_1,v_2]
=
[v_1,v_2]-[v_0,v_2]+[v_0,v_1].
\]

So the boundary of a filled triangle is a \(1\)-cycle. But it is a cycle of a special kind: it is also a boundary, so it represents zero in \(H_1\).

\subsubsection*{Geometric meaning}

In dimension \(1\):
\begin{itemize}
	\item cycles are closed loops,
	\item boundaries are loops that bound filled-in regions,
	\item homology detects \(1\)-dimensional holes.
\end{itemize}

Thus a loop is nontrivial in homology only if it closes up and does \emph{not} bound a surface in the space.

\subsubsection*{Example 1: the circle \(S^1\)}

The whole circle is a \(1\)-cycle. It does not bound any \(2\)-chain inside \(S^1\), so it survives in homology:
\[
H_1(S^1;\Z)\cong \Z.
\]

\subsubsection*{Example 2: the disk \(D^2\)}

The boundary circle of the disk is also a \(1\)-cycle, but now it bounds the interior of the disk. Therefore it is a boundary and becomes zero in homology:
\[
H_1(D^2;\Z)=0.
\]

This is one of the most important basic examples.

\subsubsection*{Example 3: the boundary of a triangle}

The boundary of a triangle is a closed loop, so it is a \(1\)-cycle. But if the triangle interior is included in the space, then that loop is also a boundary. Hence it is trivial in \(H_1\).

So:
\begin{itemize}
	\item triangle boundary alone \(\simeq S^1\): nontrivial \(H_1\),
	\item filled triangle \(\simeq D^2\): trivial \(H_1\).
\end{itemize}

\subsection{Dimension \(2\): shells and surfaces}

A \(2\)-chain is a formal integer sum of oriented triangles:
\[
c=n_1\sigma_1+n_2\sigma_2+\cdots+n_r\sigma_r.
\]

A \(2\)-cycle is a \(2\)-chain with zero boundary:
\[
\partial_2(c)=0.
\]

This means that every edge appearing in the boundary of one triangle is canceled by the opposite occurrence of that same edge from another triangle.

\subsubsection*{Picture intuition}

A \(2\)-cycle is a closed shell made out of triangles.

Typical examples are:
\begin{itemize}
	\item the surface of a tetrahedron,
	\item the boundary of a cube after triangulation,
	\item any triangulated sphere.
\end{itemize}

\subsubsection*{What is a \(2\)-boundary?}

A \(2\)-boundary is the boundary of a \(3\)-chain. So it comes from the boundary of a \(3\)-dimensional solid built from tetrahedra or other \(3\)-dimensional cells.

\subsubsection*{Geometric meaning}

In dimension \(2\):
\begin{itemize}
	\item cycles are closed shells,
	\item boundaries are shells that bound solid \(3\)-dimensional regions,
	\item homology detects \(2\)-dimensional voids.
\end{itemize}

Thus a shell is nontrivial in homology only if it closes up and does \emph{not} bound a solid inside the space.

\subsubsection*{Example 1: the sphere \(S^2\)}

The whole sphere is a \(2\)-cycle. In a triangulation, one sums all faces with consistent orientation, and every edge cancels. So the boundary is zero.

But inside \(S^2\) itself there is no \(3\)-chain filling it. Hence it is not a boundary in \(S^2\), so it survives in homology:
\[
H_2(S^2;\Z)\cong \Z.
\]

\subsubsection*{Example 2: the ball \(D^3\)}

The boundary sphere of the ball is again a \(2\)-cycle, but now it is the boundary of the whole solid ball. Hence it is trivial in homology:
\[
H_2(D^3;\Z)=0.
\]

This is the exact analogue of the circle-versus-disk example one dimension lower.

\subsubsection*{Example 3: the boundary of a tetrahedron}

The boundary of a tetrahedron consists of four triangular faces. With consistent orientation, their sum is a \(2\)-cycle, because every edge appears twice with opposite signs.

If we look only at the surface of the tetrahedron, this gives the generator of \(H_2(S^2;\Z)\). If instead we include the solid tetrahedron, then that same \(2\)-cycle becomes a boundary.

\subsection{The pattern in dimensions \(0\), \(1\), and \(2\)}

There is a perfect pattern:

\subsubsection*{Dimension \(0\)}
\begin{itemize}
	\item cycles: points,
	\item boundaries: differences of endpoints of edges,
	\item homology detects connected components.
\end{itemize}

\subsubsection*{Dimension \(1\)}
\begin{itemize}
	\item cycles: closed loops,
	\item boundaries: boundaries of filled surfaces,
	\item homology detects \(1\)-dimensional holes.
\end{itemize}

\subsubsection*{Dimension \(2\)}
\begin{itemize}
	\item cycles: closed shells,
	\item boundaries: boundaries of solids,
	\item homology detects \(2\)-dimensional voids.
\end{itemize}

\subsection{A slogan}

A cycle is \emph{closed}.

A boundary is \emph{closed for a trivial reason}, because it comes from the boundary of something higher-dimensional.

Homology records the cycles that remain after boundaries are declared trivial.

That is:
\[
H_n(X)=Z_n(X)/B_n(X).
\]

\subsection{Three fundamental comparisons}

These are among the most important examples to remember.

\subsubsection*{Circle versus disk}

The boundary circle is a \(1\)-cycle.

\begin{itemize}
	\item In \(S^1\), it is not a boundary, so it survives:
	\[
	H_1(S^1;\Z)\cong \Z.
	\]
	\item In \(D^2\), it bounds the disk, so it becomes trivial:
	\[
	H_1(D^2;\Z)=0.
	\]
\end{itemize}

\subsubsection*{Sphere versus ball}

The boundary sphere is a \(2\)-cycle.

\begin{itemize}
	\item In \(S^2\), it is not a boundary, so it survives:
	\[
	H_2(S^2;\Z)\cong \Z.
	\]
	\item In \(D^3\), it bounds the ball, so it becomes trivial:
	\[
	H_2(D^3;\Z)=0.
	\]
\end{itemize}

\subsubsection*{Connected points}

A point is always a \(0\)-cycle.

\begin{itemize}
	\item If two points are connected by a path, their difference is a boundary.
	\item Therefore all points in one connected component represent the same class.
\end{itemize}

So for connected \(X\),
\[
H_0(X;\Z)\cong \Z.
\]

\subsection{A compact geometric table}

\[
\begin{array}{c|c|c|c}
	\text{Dimension} & \text{Cycle} & \text{Boundary} & \text{What homology detects} \\ \hline
	0 & \text{points} & \text{endpoints of edges} & \text{connected components} \\
	1 & \text{closed loops} & \text{boundaries of surfaces} & \text{\(1\)-dimensional holes} \\
	2 & \text{closed shells} & \text{boundaries of solids} & \text{\(2\)-dimensional voids}
\end{array}
\]

\subsection{Best mental picture}

It is useful to remember homology in low dimensions this way:

\begin{itemize}
	\item \(H_0\): how many separate pieces does the space have?
	\item \(H_1\): which loops do not bound a surface?
	\item \(H_2\): which shells do not bound a solid?
\end{itemize}

This is the geometric meaning of cycles and boundaries in dimensions \(0\), \(1\), and \(2\).

\subsection{Summary}

The low-dimensional picture of homology can be summarized as follows:
\begin{itemize}
	\item In dimension \(0\), all points are cycles, and homology identifies points lying in the same connected component.
	\item In dimension \(1\), closed loops are cycles, but loops bounding filled-in regions are trivial in homology.
	\item In dimension \(2\), closed shells are cycles, but shells bounding solids are trivial in homology.
\end{itemize}

Thus homology is a way of detecting genuine holes by separating cycles from those cycles that arise as boundaries.

\section{\texorpdfstring{$k$}{k}-simplices and simplicial homology}

\subsection{Introduction}

Simplicial homology is one of the cleanest and most geometric ways to define homology groups. The basic idea is simple: first break a space into elementary pieces called simplices, then record how these pieces fit together by algebraic formulas, and finally extract from those formulas the groups that measure holes.

The elementary building blocks are the \emph{\(k\)-simplices}. They are the simplest geometric objects in each dimension:
\begin{itemize}
	\item a \(0\)-simplex is a point,
	\item a \(1\)-simplex is a line segment,
	\item a \(2\)-simplex is a filled triangle,
	\item a \(3\)-simplex is a filled tetrahedron,
	\item and for higher \(k\), a \(k\)-simplex is the higher-dimensional analogue of these.
\end{itemize}

A space built from simplices in a compatible way is called a \emph{simplicial complex}. Once such a complex is given, one defines chain groups, boundary maps, cycles, boundaries, and finally the homology groups
\[
H_k(K)=Z_k(K)/B_k(K).
\]
These groups measure the \(k\)-dimensional holes of the space.

The purpose of this section is to present carefully:
\begin{itemize}
	\item the definition of a \(k\)-simplex,
	\item the notion of faces and orientation,
	\item the definition of a simplicial complex,
	\item the construction of simplicial chain groups and boundary operators,
	\item the definition of simplicial homology,
	\item the main differences between the key notions,
	\item and many examples.
\end{itemize}

\subsection{Geometric \texorpdfstring{$k$}{k}-simplices}

\subsubsection{Affine independence}

Before defining a simplex, one needs the notion of affine independence.

Points \(v_0,\dots,v_k\in \mathbb{R}^n\) are said to be \emph{affinely independent} if the vectors
\[
v_1-v_0,\quad v_2-v_0,\quad \dots,\quad v_k-v_0
\]
are linearly independent in \(\mathbb{R}^n\).

Geometrically, this means that the points do not all lie in a smaller affine subspace than expected. For example:
\begin{itemize}
	\item two distinct points are affinely independent,
	\item three points in the plane are affinely independent if and only if they are not collinear,
	\item four points in \(\mathbb{R}^3\) are affinely independent if and only if they do not all lie in one plane.
\end{itemize}

\subsubsection{Definition of a geometric simplex}

Let \(v_0,\dots,v_k\in \mathbb{R}^n\) be affinely independent. The \emph{\(k\)-simplex} spanned by these vertices is the set
\[
[v_0,\dots,v_k]
=
\left\{
\sum_{i=0}^k t_i v_i
\;\middle|\;
t_i\ge 0,\ \sum_{i=0}^k t_i=1
\right\}.
\]
Thus a simplex is the set of all convex combinations of its vertices.

This is the natural higher-dimensional generalization of a segment, a triangle, and a tetrahedron.

\subsubsection{Examples of simplices}

\paragraph{A \(0\)-simplex.}
A \(0\)-simplex is simply one point:
\[
[v_0].
\]
It has dimension \(0\).

\paragraph{A \(1\)-simplex.}
A \(1\)-simplex is the line segment joining two distinct points:
\[
[v_0,v_1].
\]
It has dimension \(1\).

\paragraph{A \(2\)-simplex.}
A \(2\)-simplex is a filled triangle determined by three non-collinear points:
\[
[v_0,v_1,v_2].
\]
It includes the interior, the edges, and the vertices.

\paragraph{A \(3\)-simplex.}
A \(3\)-simplex is a filled tetrahedron:
\[
[v_0,v_1,v_2,v_3].
\]
It includes its interior, triangular faces, edges, and vertices.

\paragraph{A \(4\)-simplex.}
A \(4\)-simplex is the four-dimensional analogue of a tetrahedron. It cannot be visualized directly in ordinary three-dimensional space, but its algebraic definition is exactly the same:
\[
[v_0,v_1,v_2,v_3,v_4]
=
\left\{
\sum_{i=0}^4 t_i v_i
\;\middle|\;
t_i\ge 0,\ \sum_{i=0}^4 t_i=1
\right\},
\]
where the five points are affinely independent.

\subsubsection{The standard simplex}

A very important model is the \emph{standard \(k\)-simplex}
\[
\Delta^k
=
\left\{
(t_0,\dots,t_k)\in \mathbb{R}^{k+1}
\;\middle|\;
t_i\ge 0,\ \sum_{i=0}^k t_i=1
\right\}.
\]
This is the simplex whose vertices are the standard basis vectors
\[
e_0=(1,0,\dots,0),\quad e_1=(0,1,0,\dots,0),\quad \dots,\quad e_k=(0,\dots,0,1).
\]
Every geometric \(k\)-simplex is affinely equivalent to \(\Delta^k\).

Examples:
\[
\Delta^0=\{(1)\},
\]
which is a point,
\[
\Delta^1=\{(t_0,t_1)\in \mathbb{R}^2: t_0,t_1\ge 0,\ t_0+t_1=1\},
\]
which is a segment,
\[
\Delta^2=\{(t_0,t_1,t_2)\in \mathbb{R}^3: t_i\ge 0,\ t_0+t_1+t_2=1\},
\]
which is a triangle,
and similarly \(\Delta^3\) is a tetrahedron.

\subsection{Faces of a simplex}

A face of a simplex is obtained by deleting one or more vertices.

If
\[
\sigma=[v_0,\dots,v_k],
\]
then any simplex of the form
\[
[v_{i_0},\dots,v_{i_r}]
\quad\text{with}\quad
0\le i_0<\cdots<i_r\le k
\]
is a face of \(\sigma\). Its dimension is \(r\).

In particular, the codimension-\(1\) faces are obtained by deleting exactly one vertex:
\[
[v_1,\dots,v_k],\quad
[v_0,v_2,\dots,v_k],\quad
\dots,\quad
[v_0,\dots,v_{k-1}].
\]

\subsubsection{Examples of faces}

\paragraph{Faces of a segment.}
For the \(1\)-simplex
\[
[v_0,v_1],
\]
the faces are the vertices
\[
[v_0],\qquad [v_1].
\]

\paragraph{Faces of a triangle.}
For the \(2\)-simplex
\[
[v_0,v_1,v_2],
\]
the faces are:
\begin{itemize}
	\item the three \(0\)-faces:
	\[
	[v_0],\quad [v_1],\quad [v_2],
	\]
	\item the three \(1\)-faces:
	\[
	[v_0,v_1],\quad [v_0,v_2],\quad [v_1,v_2],
	\]
	\item and the simplex itself.
\end{itemize}

\paragraph{Faces of a tetrahedron.}
For the \(3\)-simplex
\[
[v_0,v_1,v_2,v_3],
\]
the codimension-\(1\) faces are the four triangles
\[
[v_1,v_2,v_3],\qquad
[v_0,v_2,v_3],\qquad
[v_0,v_1,v_3],\qquad
[v_0,v_1,v_2].
\]
It also has six edges and four vertices.

\subsection{Orientation of simplices}

Orientation is essential in homology because the boundary operator uses signs.

An \emph{orientation} of a \(k\)-simplex is determined by ordering its vertices:
\[
(v_0,\dots,v_k).
\]
Two orderings define the same orientation if and only if they differ by an even permutation. If they differ by an odd permutation, then they determine opposite orientations.

Thus every simplex has exactly two possible orientations.

\subsubsection{Examples}

\paragraph{Orientation of a \(1\)-simplex.}
For a segment,
\[
[v_0,v_1]
\]
and
\[
[v_1,v_0]
\]
represent opposite orientations. Algebraically one writes
\[
[v_1,v_0]=-[v_0,v_1].
\]
So orienting a \(1\)-simplex means choosing a direction.

\paragraph{Orientation of a \(2\)-simplex.}
For a triangle,
\[
[v_0,v_1,v_2]
\]
and
\[
[v_1,v_0,v_2]
\]
have opposite orientations, since the permutation exchanging \(v_0\) and \(v_1\) is odd. Hence
\[
[v_1,v_0,v_2]=-[v_0,v_1,v_2].
\]

Geometrically this corresponds to choosing clockwise or counterclockwise order.

\subsection{Abstract simplices and abstract simplicial complexes}

It is often useful to distinguish geometric simplices from abstract simplices.

\subsubsection{Abstract simplex}

An \emph{abstract simplex} on a finite set of vertices is simply a finite nonempty set
\[
\{v_0,\dots,v_k\}.
\]
Its dimension is \(k\), because it has \(k+1\) vertices.

So abstractly, a simplex is just a finite set of vertices, without coordinates.

\subsubsection{Abstract simplicial complex}

An \emph{abstract simplicial complex} \(K\) is a collection of finite nonempty subsets of a vertex set such that whenever \(\sigma\in K\), every nonempty subset of \(\sigma\) also belongs to \(K\).

Thus:
\begin{itemize}
	\item simplices are finite sets of vertices,
	\item and every face of a simplex must again be in the complex.
\end{itemize}

This is the combinatorial version of a geometric simplicial complex.

\subsubsection{Geometric realization}

From an abstract simplicial complex one can build a topological space by realizing the vertices as points in some Euclidean space and taking the corresponding geometric simplices. This space is called the \emph{geometric realization} of the abstract complex.

In practice, one often moves freely between the abstract and geometric viewpoints.

\subsection{Simplicial complexes}

A \emph{geometric simplicial complex} \(K\) is a collection of simplices in some \(\mathbb{R}^n\) such that:
\begin{enumerate}
	\item every face of a simplex of \(K\) is also in \(K\),
	\item if \(\sigma,\tau\in K\), then \(\sigma\cap\tau\) is either empty or a common face of both.
\end{enumerate}

This means that simplices are allowed to meet only in a controlled way.

The underlying topological space of the complex is
\[
|K|=\bigcup_{\sigma\in K}\sigma.
\]

\subsubsection{Examples of simplicial complexes}

\paragraph{A segment.}
Take two vertices \(v_0,v_1\) and the edge \([v_0,v_1]\). Then
\[
K=\{[v_0],[v_1],[v_0,v_1]\}
\]
is a simplicial complex. Its geometric realization is a closed interval.

\paragraph{A filled triangle.}
Take vertices \(v_0,v_1,v_2\), all three edges, and the triangle:
\[
K=\{[v_0],[v_1],[v_2],[v_0,v_1],[v_0,v_2],[v_1,v_2],[v_0,v_1,v_2]\}.
\]
Its realization is a filled triangular disk.

\paragraph{A triangulated circle.}
Take three vertices \(v_0,v_1,v_2\) and the three edges
\[
[v_0,v_1],\qquad [v_1,v_2],\qquad [v_2,v_0],
\]
but do \emph{not} include the 2-simplex \([v_0,v_1,v_2]\). Then the realization is a circle.

\paragraph{A sphere as boundary of a tetrahedron.}
Take the boundary of a tetrahedron:
\[
[v_1,v_2,v_3],\quad
[v_0,v_2,v_3],\quad
[v_0,v_1,v_3],\quad
[v_0,v_1,v_2],
\]
together with all their faces. The realization is homeomorphic to \(S^2\).

\paragraph{Two disconnected points.}
Take two vertices and no edge between them. This is a simplicial complex with two connected components.

\subsection{Important differences between the main notions}

Because the terminology is similar, it is important to separate several ideas clearly.

\subsubsection{Simplex versus simplicial complex}

A \emph{simplex} is one basic geometric piece, such as:
\begin{itemize}
	\item one point,
	\item one segment,
	\item one triangle,
	\item one tetrahedron.
\end{itemize}

A \emph{simplicial complex} is a whole space made by gluing many simplices together compatibly.

So:
\begin{itemize}
	\item a simplex is one building block,
	\item a simplicial complex is an entire structure built from such blocks.
\end{itemize}

\subsubsection{Geometric simplex versus abstract simplex}

A \emph{geometric simplex} is an actual subset of Euclidean space, given by convex combinations of vertices.

An \emph{abstract simplex} is only a finite set of vertices, with no geometric coordinates.

Thus:
\begin{itemize}
	\item geometric simplices live inside \(\mathbb{R}^n\),
	\item abstract simplices are purely combinatorial.
\end{itemize}

\subsubsection{Boundary of a simplex versus boundary of a space}

The \emph{boundary of a simplex} is the alternating sum of its codimension-\(1\) faces in the chain complex.

The \emph{topological boundary} of a subspace in a topological sense is a different notion. In simplicial homology, the word ``boundary'' usually means algebraic boundary.

\subsubsection{Cycle versus boundary}

A \emph{cycle} is a chain whose boundary is zero.

A \emph{boundary} is a chain that is itself the boundary of a higher-dimensional chain.

Every boundary is automatically a cycle, because the boundary of a boundary is always zero. But not every cycle is a boundary. Homology is precisely about measuring the difference.

\subsection{Simplicial chain groups}

Let \(K\) be a simplicial complex.

For each \(k\ge 0\), the group of \emph{simplicial \(k\)-chains} is the free abelian group generated by the oriented \(k\)-simplices of \(K\), subject to the relation that reversing orientation changes the sign.

This group is denoted by \(C_k(K)\).

Thus an element of \(C_k(K)\) is a finite formal sum
\[
c=\sum_i n_i \sigma_i,
\qquad n_i\in \mathbb{Z},
\]
where each \(\sigma_i\) is an oriented \(k\)-simplex.

\subsubsection{Examples of chain groups}

\paragraph{A single point.}
If \(K\) consists of one vertex \(v_0\), then
\[
C_0(K)\cong \mathbb{Z},
\qquad
C_k(K)=0 \quad\text{for } k\ge 1.
\]

\paragraph{A segment.}
If \(K\) consists of vertices \(v_0,v_1\) and one edge \([v_0,v_1]\), then
\[
C_0(K)\cong \mathbb{Z}\langle [v_0],[v_1]\rangle \cong \mathbb{Z}^2,
\]
\[
C_1(K)\cong \mathbb{Z}\langle [v_0,v_1]\rangle \cong \mathbb{Z},
\]
and \(C_k(K)=0\) for \(k\ge 2\).

\paragraph{A filled triangle.}
For the filled triangle,
\[
C_0(K)\cong \mathbb{Z}^3,\qquad
C_1(K)\cong \mathbb{Z}^3,\qquad
C_2(K)\cong \mathbb{Z},
\]
with bases given by the vertices, edges, and the single 2-simplex.

\subsubsection{Concrete examples of chains}

If the complex is the boundary of a triangle, then a \(1\)-chain might be
\[
[v_0,v_1]+[v_1,v_2]-[v_0,v_2].
\]
This is an algebraic combination of edges.

In the filled triangle, a \(2\)-chain is simply
\[
[v_0,v_1,v_2].
\]

It is important to remember that chains are \emph{formal sums}. They are not just subsets of the space; they are algebraic objects.

\subsection{The boundary operator}

The key map in simplicial homology is the boundary operator
\[
\partial_k:C_k(K)\to C_{k-1}(K).
\]

On an oriented \(k\)-simplex it is defined by
\[
\partial_k[v_0,\dots,v_k]
=
\sum_{i=0}^k (-1)^i [v_0,\dots,\widehat{v_i},\dots,v_k],
\]
where \(\widehat{v_i}\) means that \(v_i\) is omitted.

One then extends this map linearly to all chains.

\subsubsection{Examples of the boundary map}

\paragraph{Boundary of a \(0\)-simplex.}
A point has no lower-dimensional faces, so
\[
\partial_0[v_0]=0.
\]

\paragraph{Boundary of a \(1\)-simplex.}
For a segment,
\[
\partial_1[v_0,v_1]
=
[v_1]-[v_0].
\]
Thus the boundary of an oriented edge is its endpoint minus its starting point.

\paragraph{Boundary of a \(2\)-simplex.}
For a triangle,
\[
\partial_2[v_0,v_1,v_2]
=
[v_1,v_2]-[v_0,v_2]+[v_0,v_1].
\]
So the boundary of a triangle is the alternating sum of its three edges.

\paragraph{Boundary of a \(3\)-simplex.}
For a tetrahedron,
\[
\partial_3[v_0,v_1,v_2,v_3]
=
[v_1,v_2,v_3]
-
[v_0,v_2,v_3]
+
[v_0,v_1,v_3]
-
[v_0,v_1,v_2].
\]

\subsubsection{Why the signs matter}

Without the alternating signs, the fundamental identity \(\partial^2=0\) would fail. The signs guarantee the pairwise cancellation of faces when the boundary is taken twice.

\subsection{The identity \texorpdfstring{$\partial^2=0$}{d squared = 0}}

A central theorem of homology theory is:
\[
\partial_{k-1}\circ \partial_k = 0
\qquad \text{for all } k\ge 1.
\]
Equivalently,
\[
\partial^2=0.
\]

This means that the boundary of a boundary is always zero.

\subsubsection{Example: a triangle}

Take
\[
\sigma=[v_0,v_1,v_2].
\]
Then
\[
\partial \sigma
=
[v_1,v_2]-[v_0,v_2]+[v_0,v_1].
\]
Apply \(\partial\) again:
\[
\partial\partial \sigma
=
([v_2]-[v_1]) - ([v_2]-[v_0]) + ([v_1]-[v_0]).
\]
Now simplify:
\[
[v_2]-[v_1]-[v_2]+[v_0]+[v_1]-[v_0]=0.
\]
Everything cancels.

Geometrically, the edges of the boundary of a filled triangle have no further boundary left over.

\subsubsection{Example: a tetrahedron}

Similarly, when taking the boundary of a tetrahedron and then the boundary again, every edge appears twice with opposite signs, so the total sum is zero.

\subsection{Cycles and boundaries}

Now the two fundamental subgroups can be defined.

\subsubsection{Cycles}

A \(k\)-chain \(c\in C_k(K)\) is called a \emph{\(k\)-cycle} if
\[
\partial_k(c)=0.
\]
The group of \(k\)-cycles is denoted
\[
Z_k(K)=\ker(\partial_k).
\]

Thus a cycle is a chain with no boundary.

\subsubsection{Boundaries}

A \(k\)-chain \(c\in C_k(K)\) is called a \emph{\(k\)-boundary} if there exists a \((k+1)\)-chain \(d\in C_{k+1}(K)\) such that
\[
c=\partial_{k+1}(d).
\]
The group of \(k\)-boundaries is denoted
\[
B_k(K)=\operatorname{im}(\partial_{k+1}).
\]

Thus a boundary is a chain that comes from the boundary of something one dimension higher.

\subsubsection{Relation between them}

Since \(\partial^2=0\), every boundary is a cycle:
\[
B_k(K)\subseteq Z_k(K).
\]
Indeed, if \(c=\partial d\), then
\[
\partial c=\partial(\partial d)=0.
\]

This inclusion is exactly what makes the quotient \(Z_k/B_k\) meaningful.

\subsection{Simplicial homology groups}

The \emph{\(k\)-th simplicial homology group} of \(K\) is defined to be
\[
H_k(K)=Z_k(K)/B_k(K).
\]

This means:
\begin{itemize}
	\item start with all \(k\)-cycles,
	\item then identify two cycles if their difference is a boundary.
\end{itemize}

So homology classes are cycles modulo those that are trivial because they bound something.

\subsubsection{Geometric meaning}

The group \(H_k(K)\) measures \(k\)-dimensional holes:
\begin{itemize}
	\item \(H_0\) measures connected components,
	\item \(H_1\) measures one-dimensional holes, such as essential loops,
	\item \(H_2\) measures two-dimensional cavities or shells,
	\item and so on.
\end{itemize}

\subsection{Intuition in low dimensions}

\subsubsection{\texorpdfstring{$H_0$}{H0}}

The group \(H_0\) measures connected components. If a space has \(m\) connected components, then typically
\[
H_0\cong \mathbb{Z}^m.
\]

\subsubsection{\texorpdfstring{$H_1$}{H1}}

The group \(H_1\) measures independent loops that do not bound any filled region in the complex.

Examples:
\begin{itemize}
	\item the circle has one essential loop,
	\item the torus has two independent loops,
	\item a filled disk has no essential loop, because its boundary bounds the disk.
\end{itemize}

\subsubsection{\texorpdfstring{$H_2$}{H2}}

The group \(H_2\) measures closed surfaces that do not bound a 3-chain inside the complex.

For example:
\begin{itemize}
	\item the sphere has one fundamental \(2\)-cycle,
	\item the torus also has one fundamental \(2\)-cycle,
	\item the circle has no \(2\)-dimensional homology.
\end{itemize}

\subsection{Many detailed examples}

\subsubsection{Example 1: a single point}

Let \(K=\{[v_0]\}\).

Then
\[
C_0(K)\cong \mathbb{Z},\qquad C_k(K)=0 \text{ for } k\ge 1.
\]
Since \(\partial_0=0\), every \(0\)-chain is a \(0\)-cycle:
\[
Z_0(K)=C_0(K)\cong \mathbb{Z}.
\]
Also
\[
B_0(K)=\operatorname{im}(\partial_1)=0,
\]
because \(C_1(K)=0\).

Hence
\[
H_0(K)\cong \mathbb{Z}.
\]
And for \(k\ge 1\),
\[
H_k(K)=0.
\]

So a point has one connected component and no higher-dimensional holes.

\subsubsection{Example 2: two disconnected points}

Let \(K\) have vertices \(v_0,v_1\) and no edges.

Then
\[
C_0(K)\cong \mathbb{Z}^2,\qquad C_k(K)=0 \text{ for } k\ge 1.
\]
Again
\[
Z_0(K)=C_0(K)\cong \mathbb{Z}^2,
\qquad
B_0(K)=0.
\]
So
\[
H_0(K)\cong \mathbb{Z}^2.
\]
This reflects the fact that there are two connected components.

\subsubsection{Example 3: a segment}

Let \(K\) have vertices \(v_0,v_1\) and one edge \([v_0,v_1]\).

Then
\[
C_1(K)\cong \mathbb{Z}\langle [v_0,v_1]\rangle,
\qquad
C_0(K)\cong \mathbb{Z}\langle [v_0],[v_1]\rangle.
\]
The boundary map is
\[
\partial_1[v_0,v_1]=[v_1]-[v_0].
\]

To compute \(H_1\), note that a \(1\)-chain is
\[
n[v_0,v_1].
\]
Its boundary is
\[
\partial_1(n[v_0,v_1])=n([v_1]-[v_0]).
\]
This is zero only if \(n=0\), so
\[
Z_1(K)=0.
\]
Hence
\[
H_1(K)=0.
\]

For \(H_0\), the image of \(\partial_1\) is generated by
\[
[v_1]-[v_0].
\]
Thus
\[
H_0(K)
=
\frac{\mathbb{Z}\langle [v_0],[v_1]\rangle}
{\mathbb{Z}\langle [v_1]-[v_0]\rangle}
\cong \mathbb{Z}.
\]
So the segment is connected and has no \(1\)-dimensional hole.

\subsubsection{Example 4: the boundary of a triangle, i.e. a circle}

Let \(K\) have vertices \(v_0,v_1,v_2\) and edges
\[
[v_0,v_1],\qquad [v_1,v_2],\qquad [v_2,v_0],
\]
but no \(2\)-simplex.

Then \(K\) is a triangulation of \(S^1\).

The chain group \(C_1(K)\) is free abelian on the three oriented edges, and \(C_2(K)=0\).

Consider the \(1\)-chain
\[
c=[v_0,v_1]+[v_1,v_2]+[v_2,v_0].
\]
Its boundary is
\[
\partial c
=
([v_1]-[v_0]) + ([v_2]-[v_1]) + ([v_0]-[v_2]) = 0.
\]
So \(c\) is a \(1\)-cycle.

Since \(C_2(K)=0\), there are no nonzero \(1\)-boundaries:
\[
B_1(K)=0.
\]
Therefore the class of \(c\) is nontrivial in homology, and
\[
H_1(K)\cong \mathbb{Z}.
\]

Also the space is connected, so
\[
H_0(K)\cong \mathbb{Z}.
\]

Thus
\[
H_0(S^1)\cong \mathbb{Z},
\qquad
H_1(S^1)\cong \mathbb{Z},
\qquad
H_k(S^1)=0 \text{ for } k\ge 2.
\]

\subsubsection{Example 5: a filled triangle}

Now let \(K\) contain the triangle \([v_0,v_1,v_2]\) itself, as well as all its faces.

Then
\[
C_2(K)\cong \mathbb{Z}\langle [v_0,v_1,v_2]\rangle.
\]
The boundary of this \(2\)-simplex is
\[
\partial_2[v_0,v_1,v_2]
=
[v_1,v_2]-[v_0,v_2]+[v_0,v_1].
\]
This is exactly the triangular loop, with compatible orientation.

Hence the loop that was nontrivial in the previous example is now a boundary. Therefore
\[
H_1(K)=0.
\]

The space is connected, so
\[
H_0(K)\cong \mathbb{Z}.
\]

There is no nontrivial \(2\)-cycle, because the single \(2\)-simplex has nonzero boundary. So
\[
H_2(K)=0.
\]

Thus the filled triangle has the same homology as a point:
\[
H_0(K)\cong \mathbb{Z},
\qquad
H_k(K)=0 \text{ for } k\ge 1.
\]

This example is extremely important: it shows the difference between a loop and a loop that bounds a filled region.

\subsubsection{Example 6: the boundary of a tetrahedron, i.e. a sphere}

Take the four triangular faces of a tetrahedron, together with all edges and vertices, but not the solid interior. The resulting simplicial complex is homeomorphic to \(S^2\).

Let the four oriented faces be
\[
\sigma_1,\sigma_2,\sigma_3,\sigma_4,
\]
chosen consistently. Then the sum
\[
c=\sigma_1+\sigma_2+\sigma_3+\sigma_4
\]
is a \(2\)-cycle, because every edge appears twice with opposite orientations and cancels.

Since there is no \(3\)-simplex in the boundary complex,
\[
B_2(K)=0.
\]
Therefore
\[
H_2(K)\cong \mathbb{Z}.
\]

Also \(S^2\) has no essential loops, so
\[
H_1(K)=0.
\]
And the space is connected, so
\[
H_0(K)\cong \mathbb{Z}.
\]

Hence
\[
H_0(S^2)\cong \mathbb{Z},\qquad
H_1(S^2)=0,\qquad
H_2(S^2)\cong \mathbb{Z}.
\]

\subsubsection{Example 7: two disjoint circles}

If a simplicial complex is the disjoint union of two triangulated circles, then each circle contributes one connected component and one essential \(1\)-cycle. Hence
\[
H_0\cong \mathbb{Z}^2,\qquad H_1\cong \mathbb{Z}^2.
\]

\subsubsection{Example 8: a torus}

A torus can be triangulated by a finite simplicial complex. Its simplicial homology is
\[
H_0(T^2)\cong \mathbb{Z},\qquad
H_1(T^2)\cong \mathbb{Z}^2,\qquad
H_2(T^2)\cong \mathbb{Z}.
\]

Geometrically:
\begin{itemize}
	\item \(H_0\cong \mathbb{Z}\) because the torus is connected,
	\item \(H_1\cong \mathbb{Z}^2\) because there are two independent fundamental loops: the meridian and the longitude,
	\item \(H_2\cong \mathbb{Z}\) because the whole torus surface gives a fundamental \(2\)-cycle.
\end{itemize}

\subsubsection{Example 9: a hollow square and a filled square}

A square can be triangulated in two different ways giving different homology.

\paragraph{Hollow square.}
Take the boundary edges only, triangulated as a polygonal loop. The resulting space is homeomorphic to \(S^1\), so
\[
H_1\cong \mathbb{Z}.
\]

\paragraph{Filled square.}
Now fill the square by dividing it into two triangles. Then the boundary loop becomes the boundary of the sum of those two triangles, so
\[
H_1=0.
\]
This is another illustration of how filling in a loop kills a homology class.

\subsection{Cycle versus boundary: the central distinction}

This is the conceptual heart of homology.

\subsubsection{A cycle}

A cycle is something closed.

Examples:
\begin{itemize}
	\item a closed polygonal loop in a circle,
	\item the entire surface of a sphere,
	\item a sum of edges that balances so that every vertex cancels.
\end{itemize}

\subsubsection{A boundary}

A boundary is something closed for a trivial reason: it is the boundary of a higher-dimensional piece.

Examples:
\begin{itemize}
	\item the boundary of a filled triangle,
	\item the boundary of a square filled by two triangles,
	\item the boundary of a solid tetrahedron.
\end{itemize}

\subsubsection{Why homology uses the quotient \texorpdfstring{$Z_k/B_k$}{Zk/Bk}}

Homology wants to keep the genuinely interesting closed objects and ignore the trivial ones. That is why one divides cycles by boundaries:
\[
H_k=Z_k/B_k.
\]

If a cycle bounds something, then in homology it is declared to be zero.

If a cycle does \emph{not} bound anything, then it represents a nontrivial homology class and therefore detects a hole.

\subsection{Simplicial homology versus singular homology}

Another important difference is the following.

\subsubsection{Simplicial homology}

Simplicial homology is defined for a simplicial complex \(K\), using its simplices directly.

It is concrete, combinatorial, and often easy to compute.

\subsubsection{Singular homology}

Singular homology is defined for an arbitrary topological space \(X\), using all continuous maps
\[
\sigma:\Delta^k\to X.
\]
These maps are called singular simplices.

This theory works for all spaces, not only triangulated ones.

\subsubsection{Relation between them}

If a space admits a triangulation, then simplicial homology and singular homology give the same homology groups. Thus simplicial homology is a very concrete model of the general theory.

\subsection{Why simplicial homology is useful}

Simplicial homology is important for several reasons.

\begin{enumerate}
	\item It converts topology into algebra. A space becomes a chain complex of free abelian groups.
	\item It makes holes computable. One can write matrices for boundary maps and compute kernels and images explicitly.
	\item It is geometrically intuitive. Vertices, edges, triangles, and tetrahedra are easy to visualize.
	\item It serves as a model for more general homology theories.
	\item It is foundational in algebraic topology, combinatorial topology, and computational topology.
\end{enumerate}

\subsection{A compact summary of the basic definitions}

For quick reference, here are the main definitions again.

\paragraph{A \(k\)-simplex.}
If \(v_0,\dots,v_k\) are affinely independent, then
\[
[v_0,\dots,v_k]
=
\left\{
\sum_{i=0}^k t_i v_i
\;\middle|\;
t_i\ge 0,\ \sum t_i=1
\right\}.
\]

\paragraph{A simplicial complex.}
A collection of simplices closed under faces, such that two simplices intersect only along a common face.

\paragraph{The chain group.}
\[
C_k(K)=\text{free abelian group on oriented \(k\)-simplices of \(K\)}.
\]

\paragraph{The boundary operator.}
\[
\partial_k[v_0,\dots,v_k]
=
\sum_{i=0}^k (-1)^i [v_0,\dots,\widehat{v_i},\dots,v_k].
\]

\paragraph{Cycles.}
\[
Z_k(K)=\ker(\partial_k).
\]

\paragraph{Boundaries.}
\[
B_k(K)=\operatorname{im}(\partial_{k+1}).
\]

\paragraph{Homology.}
\[
H_k(K)=Z_k(K)/B_k(K).
\]

\subsection{Final summary}

A \(k\)-simplex is the simplest possible geometric object of dimension \(k\). It is built from \(k+1\) affinely independent vertices and consists of all convex combinations of those vertices. Simplices serve as the building blocks of simplicial complexes, which are spaces assembled from simplices in a compatible way.

Once a simplicial complex is given, one forms chain groups from oriented simplices, defines the boundary operator by alternating sums of faces, proves that \(\partial^2=0\), and then defines cycles, boundaries, and homology groups. The resulting simplicial homology groups detect connected components, loops, cavities, and higher-dimensional holes.

The main conceptual distinctions are:
\begin{itemize}
	\item simplex versus simplicial complex,
	\item geometric versus abstract simplex,
	\item cycle versus boundary,
	\item simplicial homology versus singular homology.
\end{itemize}

The main geometric message is this: homology studies closed objects up to those that bound something. In this way, simplicial homology turns the combinatorics of a triangulation into algebraic invariants of the underlying space.

\section{Orientations, chains, boundaries, and explicit computations}

\subsection{Introduction}

After defining simplices and simplicial homology abstractly, the next important step is to understand how the machinery works in practice. The key ingredients are:
\begin{itemize}
	\item orientations of simplices,
	\item simplicial chains as algebraic combinations of simplices,
	\item the boundary operator,
	\item cycles and boundaries,
	\item and explicit computations of homology groups in concrete examples.
\end{itemize}

This section develops these ideas in detail. The emphasis is on seeing exactly how the algebra reflects the geometry. In particular, we will study:
\begin{itemize}
	\item a single triangle and its boundary,
	\item the boundary of a tetrahedron, which gives a triangulation of \(S^2\),
	\item and the torus, where one sees two independent \(1\)-dimensional homology classes.
\end{itemize}

\subsection{Orientation revisited}

\subsubsection{Why orientation is needed}

If one wants to form algebraic sums of simplices and define a boundary map, one must be able to distinguish a simplex from the same simplex with reversed order of vertices. This is why orientations are introduced.

Geometrically:
\begin{itemize}
	\item an oriented edge has a direction,
	\item an oriented triangle has a chosen cyclic order,
	\item an oriented tetrahedron has an orientation determined by the order of its vertices.
\end{itemize}

Without orientation, the boundary formula would not carry signs, and the identity \(\partial^2=0\) would fail.

\subsubsection{Orientation of a \texorpdfstring{$1$}{1}-simplex}

For an edge \([v_0,v_1]\), the two possible orientations are
\[
[v_0,v_1]
\qquad\text{and}\qquad
[v_1,v_0].
\]
They are opposites of one another, and algebraically one writes
\[
[v_1,v_0] = -[v_0,v_1].
\]

Thus orienting an edge means choosing which endpoint is initial and which is terminal.

\subsubsection{Orientation of a \texorpdfstring{$2$}{2}-simplex}

For a triangle \([v_0,v_1,v_2]\), the possible orderings of the vertices split into two classes:
\begin{itemize}
	\item those differing by even permutations,
	\item those differing by odd permutations.
\end{itemize}
The two classes represent the two possible orientations.

For example,
\[
[v_0,v_1,v_2]
\]
and
\[
[v_1,v_2,v_0]
\]
have the same orientation, because the permutation \((0\,1\,2)\) is even. But
\[
[v_1,v_0,v_2]
=
-[v_0,v_1,v_2].
\]

Geometrically, one may think of the two orientations as clockwise and counterclockwise.

\subsubsection{Orientation of a \texorpdfstring{$3$}{3}-simplex}

For a tetrahedron \([v_0,v_1,v_2,v_3]\), the orientation is again determined by the order of the vertices up to even permutation. Swapping two vertices reverses the orientation.

In higher dimensions the same rule continues: the orientation of a \(k\)-simplex is given by an ordering of its \(k+1\) vertices, modulo even permutations.

\subsection{Simplicial chains in detail}

Let \(K\) be a simplicial complex.

\subsubsection{Definition of the chain groups}

For each \(k\ge 0\), the group \(C_k(K)\) of simplicial \(k\)-chains is the free abelian group generated by the oriented \(k\)-simplices of \(K\), with the relation that reversing orientation multiplies by \(-1\).

Thus every \(k\)-chain has the form
\[
c = \sum_{i=1}^N n_i \sigma_i,
\qquad n_i\in \mathbb{Z},
\]
where each \(\sigma_i\) is an oriented \(k\)-simplex.

The coefficients may be positive, negative, or zero.

\subsubsection{Geometric meaning of a chain}

A chain should not be confused with an actual subset of the space. It is an algebraic object.

For example, if
\[
c = [v_0,v_1] + [v_1,v_2] - [v_0,v_2],
\]
then \(c\) is not simply the union of three edges. Rather, it is a formal sum with coefficients and signs.

The signs matter because when boundaries are computed, terms cancel according to orientation.

\subsubsection{Examples of chains}

\paragraph{In dimension \(0\).}
A \(0\)-chain is a finite formal sum of vertices:
\[
2[v_0] - 3[v_1] + [v_2].
\]
So \(0\)-chains are formal integer combinations of points.

\paragraph{In dimension \(1\).}
A \(1\)-chain is a formal sum of oriented edges:
\[
[v_0,v_1] + [v_1,v_2] - 2[v_2,v_3].
\]

\paragraph{In dimension \(2\).}
A \(2\)-chain is a formal sum of oriented triangles:
\[
[v_0,v_1,v_2] - [v_0,v_2,v_3].
\]

\subsubsection{Chain groups as free abelian groups}

If a complex has finitely many \(k\)-simplices, then \(C_k(K)\) is isomorphic to \(\mathbb{Z}^m\), where \(m\) is the number of oriented \(k\)-simplices in a chosen basis.

This makes simplicial homology computable by linear algebra over \(\mathbb{Z}\).

\subsection{The boundary map in detail}

\subsubsection{Definition}

The boundary map
\[
\partial_k : C_k(K)\to C_{k-1}(K)
\]
is defined on a basis simplex by
\[
\partial_k[v_0,\dots,v_k]
=
\sum_{i=0}^k (-1)^i [v_0,\dots,\widehat{v_i},\dots,v_k],
\]
and then extended linearly.

So the boundary of a simplex is the alternating sum of its codimension-\(1\) faces.

\subsubsection{Boundary of a point}

Since a point has no lower-dimensional faces,
\[
\partial_0[v_0]=0.
\]

\subsubsection{Boundary of an edge}

For an oriented edge,
\[
\partial_1[v_0,v_1]=[v_1]-[v_0].
\]
Thus the boundary is terminal vertex minus initial vertex.

This matches the intuition from directed graphs.

\subsubsection{Boundary of a triangle}

For a triangle,
\[
\partial_2[v_0,v_1,v_2]
=
[v_1,v_2]-[v_0,v_2]+[v_0,v_1].
\]

This can also be written as
\[
\partial_2[v_0,v_1,v_2]
=
[v_0,v_1]+[v_1,v_2]+[v_2,v_0],
\]
provided one remembers that
\[
[v_2,v_0]=-[v_0,v_2].
\]

Geometrically, the boundary of the filled triangle is its oriented edge loop.

\subsubsection{Boundary of a tetrahedron}

For a tetrahedron,
\[
\partial_3[v_0,v_1,v_2,v_3]
=
[v_1,v_2,v_3]
-
[v_0,v_2,v_3]
+
[v_0,v_1,v_3]
-
[v_0,v_1,v_2].
\]
Thus the boundary of a tetrahedron is the alternating sum of its four triangular faces.

\subsection{Why \texorpdfstring{$\partial^2=0$}{d squared = 0}: a concrete explanation}

The theorem
\[
\partial_{k-1}\partial_k = 0
\]
means that every codimension-\(2\) face appears twice in the second boundary, once with a plus sign and once with a minus sign.

\subsubsection{Direct verification for a triangle}

Take
\[
\sigma=[v_0,v_1,v_2].
\]
Then
\[
\partial_2\sigma
=
[v_1,v_2]-[v_0,v_2]+[v_0,v_1].
\]
Applying \(\partial_1\), we get
\[
\partial_1\partial_2\sigma
=
([v_2]-[v_1]) - ([v_2]-[v_0]) + ([v_1]-[v_0]).
\]
Expanding:
\[
[v_2]-[v_1]-[v_2]+[v_0]+[v_1]-[v_0]=0.
\]

So the boundary of the boundary is zero.

\subsubsection{Direct verification for a tetrahedron}

Take
\[
\tau=[v_0,v_1,v_2,v_3].
\]
Then
\[
\partial_3\tau
=
[v_1,v_2,v_3]
-
[v_0,v_2,v_3]
+
[v_0,v_1,v_3]
-
[v_0,v_1,v_2].
\]
Now take the boundary of each triangle. Each edge of the tetrahedron appears exactly twice, with opposite signs, so everything cancels.

Geometrically, the surface of a tetrahedron has no boundary.

\subsection{Cycles and boundaries again, now concretely}

\subsubsection{Cycles}

A \(k\)-chain \(c\) is a cycle if
\[
\partial_k(c)=0.
\]
This means that all lower-dimensional boundary contributions cancel.

Examples:
\begin{itemize}
	\item a closed loop of edges is a \(1\)-cycle,
	\item the full surface of a tetrahedron is a \(2\)-cycle.
\end{itemize}

\subsubsection{Boundaries}

A \(k\)-chain \(c\) is a boundary if there exists a \((k+1)\)-chain \(d\) such that
\[
c=\partial_{k+1}(d).
\]

Examples:
\begin{itemize}
	\item the edge loop around a filled triangle is a \(1\)-boundary,
	\item the surface of a solid tetrahedron is a \(2\)-boundary.
\end{itemize}

\subsubsection{Key principle}

A cycle is any closed object.

A boundary is a closed object that is closed for a trivial reason: it is the boundary of something higher-dimensional.

Homology keeps track of cycles that are \emph{not} boundaries.

\subsection{Explicit computation for a triangle}

We now work through the standard triangle example in full detail.

\subsubsection{The simplicial complex}

Let \(K\) be the filled triangle with vertices \(v_0,v_1,v_2\), edges
\[
[v_0,v_1],\qquad [v_1,v_2],\qquad [v_0,v_2],
\]
and one \(2\)-simplex
\[
[v_0,v_1,v_2].
\]

\subsubsection{Chain groups}

The chain groups are:
\[
C_2(K)\cong \mathbb{Z},
\qquad
C_1(K)\cong \mathbb{Z}^3,
\qquad
C_0(K)\cong \mathbb{Z}^3.
\]

A basis may be chosen as:
\[
C_2(K)=\mathbb{Z}\langle [v_0,v_1,v_2]\rangle,
\]
\[
C_1(K)=\mathbb{Z}\langle [v_0,v_1],[v_1,v_2],[v_0,v_2]\rangle,
\]
\[
C_0(K)=\mathbb{Z}\langle [v_0],[v_1],[v_2]\rangle.
\]

\subsubsection{Boundary maps}

The boundary map in dimension \(2\) is
\[
\partial_2[v_0,v_1,v_2]
=
[v_1,v_2]-[v_0,v_2]+[v_0,v_1].
\]

The boundary map in dimension \(1\) is determined by
\[
\partial_1[v_0,v_1]=[v_1]-[v_0],
\]
\[
\partial_1[v_1,v_2]=[v_2]-[v_1],
\]
\[
\partial_1[v_0,v_2]=[v_2]-[v_0].
\]

\subsubsection{Computing \texorpdfstring{$H_2$}{H2}}

A \(2\)-chain has the form
\[
n[v_0,v_1,v_2].
\]
Its boundary is
\[
\partial_2(n[v_0,v_1,v_2])
=
n\bigl([v_1,v_2]-[v_0,v_2]+[v_0,v_1]\bigr).
\]
This is zero only if \(n=0\). Hence
\[
Z_2(K)=0.
\]
Since \(B_2(K)=0\) automatically, we get
\[
H_2(K)=0.
\]

\subsubsection{Computing \texorpdfstring{$H_1$}{H1}}

A general \(1\)-chain is
\[
a[v_0,v_1]+b[v_1,v_2]+c[v_0,v_2].
\]
Its boundary is
\[
a([v_1]-[v_0]) + b([v_2]-[v_1]) + c([v_2]-[v_0]).
\]
Collecting coefficients of vertices:
\[
\partial_1
=
(-a-c)[v_0] + (a-b)[v_1] + (b+c)[v_2].
\]
For this to be zero, we need
\[
-a-c=0,\qquad a-b=0,\qquad b+c=0.
\]
These imply
\[
a=b,\qquad c=-a.
\]
So every \(1\)-cycle has the form
\[
a\bigl([v_0,v_1]+[v_1,v_2]-[v_0,v_2]\bigr).
\]
Thus
\[
Z_1(K)\cong \mathbb{Z}.
\]

But
\[
\partial_2[v_0,v_1,v_2]
=
[v_1,v_2]-[v_0,v_2]+[v_0,v_1]
=
[v_0,v_1]+[v_1,v_2]-[v_0,v_2].
\]
So the generator of \(Z_1(K)\) is actually a boundary. Hence
\[
B_1(K)=Z_1(K),
\]
and therefore
\[
H_1(K)=0.
\]

\subsubsection{Computing \texorpdfstring{$H_0$}{H0}}

Since the triangle is connected,
\[
H_0(K)\cong \mathbb{Z}.
\]
One may also compute this directly:
\[
H_0(K)=\frac{Z_0(K)}{B_0(K)}.
\]
Because \(\partial_0=0\), all \(0\)-chains are cycles:
\[
Z_0(K)=C_0(K)\cong \mathbb{Z}^3.
\]
The subgroup \(B_0(K)\) is generated by differences of vertices arising from edges, so after quotienting one identifies all three vertices in homology. Thus
\[
H_0(K)\cong \mathbb{Z}.
\]

\subsubsection{Conclusion for the triangle}

The filled triangle has homology
\[
H_0(K)\cong \mathbb{Z},
\qquad
H_1(K)=0,
\qquad
H_2(K)=0.
\]
So it has the same homology as a point, as expected for a contractible space.

\subsection{Explicit computation for the boundary of a triangle}

Now remove the \(2\)-simplex but keep the three edges and three vertices. This is a triangulation of \(S^1\).

\subsubsection{Chain groups}

Now
\[
C_2(K)=0,
\qquad
C_1(K)\cong \mathbb{Z}^3,
\qquad
C_0(K)\cong \mathbb{Z}^3.
\]

\subsubsection{Computing \texorpdfstring{$H_1$}{H1}}

The cycle computation is exactly the same as before:
\[
Z_1(K)\cong \mathbb{Z},
\]
generated by
\[
[v_0,v_1]+[v_1,v_2]-[v_0,v_2].
\]
But now there are no \(2\)-simplices, so
\[
B_1(K)=0.
\]
Hence
\[
H_1(K)\cong \mathbb{Z}.
\]

\subsubsection{Computing \texorpdfstring{$H_0$}{H0}}

The space is connected, so
\[
H_0(K)\cong \mathbb{Z}.
\]

\subsubsection{Conclusion}

The boundary of the triangle has homology
\[
H_0(K)\cong \mathbb{Z},
\qquad
H_1(K)\cong \mathbb{Z},
\qquad
H_k(K)=0 \text{ for } k\ge 2.
\]
This is exactly the homology of the circle.

The contrast with the filled triangle is fundamental:
\begin{itemize}
	\item the loop exists in both cases,
	\item but in the filled triangle it becomes a boundary,
	\item while in the hollow triangle it remains nontrivial in homology.
\end{itemize}

\subsection{Explicit computation for the boundary of a tetrahedron}

We now turn to a standard triangulation of the sphere \(S^2\).

\subsubsection{The simplicial complex}

Let \(K\) be the boundary of a tetrahedron with vertices \(v_0,v_1,v_2,v_3\). Its \(2\)-simplices are
\[
\sigma_1=[v_1,v_2,v_3],\qquad
\sigma_2=[v_0,v_2,v_3],\qquad
\sigma_3=[v_0,v_1,v_3],\qquad
\sigma_4=[v_0,v_1,v_2].
\]

Together with all edges and vertices, these form a simplicial complex homeomorphic to \(S^2\).

\subsubsection{The fundamental \texorpdfstring{$2$}{2}-cycle}

Consider the chain
\[
c=\sigma_1-\sigma_2+\sigma_3-\sigma_4.
\]
With a consistent choice of orientations, this is the sum of all faces of the tetrahedron oriented outward.

Its boundary is zero because every edge appears twice with opposite signs. Hence
\[
c\in Z_2(K).
\]

Since the complex has no \(3\)-simplices,
\[
B_2(K)=0.
\]
Therefore \(c\) determines a nontrivial homology class, and
\[
H_2(K)\cong \mathbb{Z}.
\]

\subsubsection{Why \texorpdfstring{$H_1=0$}{H1 = 0}}

Any loop drawn on the sphere can be contracted. In the simplicial complex this means every \(1\)-cycle is the boundary of some \(2\)-chain made from the triangular faces.

Therefore
\[
H_1(K)=0.
\]

\subsubsection{Why \texorpdfstring{$H_0=\mathbb{Z}$}{H0 = Z}}

The boundary of the tetrahedron is connected, so
\[
H_0(K)\cong \mathbb{Z}.
\]

\subsubsection{Conclusion for the tetrahedral sphere}

Thus
\[
H_0(K)\cong \mathbb{Z},
\qquad
H_1(K)=0,
\qquad
H_2(K)\cong \mathbb{Z}.
\]

This is the standard homology of \(S^2\).

\subsection{Explicit picture of the sphere's \texorpdfstring{$2$}{2}-cycle}

It is worth emphasizing what the generator of \(H_2(S^2)\) looks like.

The generator is not one single triangle. A single triangle has nonzero boundary, so it is not a \(2\)-cycle.

The generator is the sum of all triangular faces of the tetrahedron, oriented consistently. The edge terms cancel pairwise, leaving zero boundary.

So the homology class in \(H_2(S^2)\) corresponds to the entire closed surface.

This is the \(2\)-dimensional analogue of how the loop around a circle represents the generator of \(H_1(S^1)\).

\subsection{The torus: geometric intuition first}

A torus \(T^2\) may be obtained from a square by identifying opposite sides:
\[
\text{left edge} \sim \text{right edge},
\qquad
\text{bottom edge} \sim \text{top edge}.
\]

This description already suggests two independent loops:
\begin{itemize}
	\item one going horizontally around the torus,
	\item one going vertically around the torus.
\end{itemize}

These give the two generators of \(H_1(T^2)\).

The whole torus surface itself gives a generator of \(H_2(T^2)\).

\subsection{A simplicial model of the torus}

To compute simplicial homology, one triangulates the square before imposing edge identifications. For example, one may divide the square into two triangles by a diagonal. After identifying opposite sides, this becomes a triangulation of the torus.

Although the full combinatorics can be written out explicitly, the main geometric conclusions are already clear from the quotient square picture.

\subsubsection{The two basic \texorpdfstring{$1$}{1}-cycles}

Let:
\begin{itemize}
	\item \(a\) be the loop that runs once horizontally,
	\item \(b\) be the loop that runs once vertically.
\end{itemize}

These are \(1\)-cycles. They do not bound \(2\)-chains on the torus, because each one winds around a genuine hole.

Thus their homology classes \([a]\) and \([b]\) are nontrivial.

Moreover, they are independent, so
\[
H_1(T^2)\cong \mathbb{Z}\oplus \mathbb{Z}.
\]

\subsubsection{The fundamental \texorpdfstring{$2$}{2}-cycle}

The torus is an oriented closed surface, so the sum of all oriented \(2\)-simplices in a triangulation gives a \(2\)-cycle.

Since the torus has no \(3\)-dimensional interior in the simplicial complex, this \(2\)-cycle is not a boundary. Hence
\[
H_2(T^2)\cong \mathbb{Z}.
\]

\subsubsection{Connectedness}

The torus is connected, so
\[
H_0(T^2)\cong \mathbb{Z}.
\]

\subsubsection{Conclusion for the torus}

Therefore
\[
H_0(T^2)\cong \mathbb{Z},
\qquad
H_1(T^2)\cong \mathbb{Z}^2,
\qquad
H_2(T^2)\cong \mathbb{Z}.
\]

\subsection{What the torus computation means geometrically}

The homology of the torus reflects three kinds of topological information:

\paragraph{One connected component.}
This gives
\[
H_0(T^2)\cong \mathbb{Z}.
\]

\paragraph{Two independent \(1\)-dimensional holes.}
These are represented by the meridian and longitude:
\[
H_1(T^2)\cong \mathbb{Z}^2.
\]

\paragraph{One \(2\)-dimensional fundamental surface.}
The whole torus surface represents
\[
H_2(T^2)\cong \mathbb{Z}.
\]

So the torus has richer homology than either the circle or the sphere:
\begin{itemize}
	\item the circle has one essential \(1\)-cycle,
	\item the sphere has one essential \(2\)-cycle,
	\item the torus has both a nontrivial \(H_1\) and a nontrivial \(H_2\).
\end{itemize}

\subsection{Comparison of the main examples}

It is useful to place the examples side by side.

\[
\begin{array}{c|c|c|c}
	\text{Space} & H_0 & H_1 & H_2 \\ \hline
	\text{point} & \mathbb{Z} & 0 & 0 \\
	\text{segment} & \mathbb{Z} & 0 & 0 \\
	S^1 & \mathbb{Z} & \mathbb{Z} & 0 \\
	\text{filled triangle} & \mathbb{Z} & 0 & 0 \\
	S^2 & \mathbb{Z} & 0 & \mathbb{Z} \\
	T^2 & \mathbb{Z} & \mathbb{Z}^2 & \mathbb{Z}
\end{array}
\]

This table shows clearly how homology detects different kinds of holes.

\subsection{Algebraic viewpoint: chain complexes}

All of the above can be organized in the language of chain complexes:
\[
\cdots \longrightarrow C_2(K)\xrightarrow{\partial_2} C_1(K)\xrightarrow{\partial_1} C_0(K)\xrightarrow{\partial_0} 0.
\]

The key condition
\[
\partial_{k-1}\partial_k=0
\]
says that this is indeed a chain complex.

Then
\[
H_k(K)=\frac{\ker(\partial_k)}{\operatorname{im}(\partial_{k+1})}.
\]

So homology is obtained by comparing:
\begin{itemize}
	\item the kernel of the boundary map in dimension \(k\),
	\item with the image of the boundary map in dimension \(k+1\).
\end{itemize}

This is the algebraic heart of simplicial homology.

\subsection{Why explicit computations matter}

The formal definition
\[
H_k(K)=Z_k(K)/B_k(K)
\]
is compact, but its meaning becomes truly clear only in examples.

The triangle and its boundary show:
\begin{itemize}
	\item how a loop may be a cycle,
	\item how filling it in makes it a boundary,
	\item and how this changes the homology group.
\end{itemize}

The tetrahedron shows:
\begin{itemize}
	\item how a closed surface becomes a \(2\)-cycle,
	\item how all edge boundaries cancel,
	\item and why the sphere has nontrivial \(H_2\).
\end{itemize}

The torus shows:
\begin{itemize}
	\item how different independent loops produce a higher-rank \(H_1\),
	\item and how a closed surface can simultaneously have nontrivial \(H_1\) and \(H_2\).
\end{itemize}

\subsection{Final summary}

Orientations allow simplices to be used algebraically with signs. Chains are formal sums of oriented simplices, and the boundary operator sends each simplex to the alternating sum of its faces. The fundamental identity \(\partial^2=0\) ensures that boundaries are always cycles.

This leads to the definitions:
\[
Z_k(K)=\ker(\partial_k),
\qquad
B_k(K)=\operatorname{im}(\partial_{k+1}),
\qquad
H_k(K)=Z_k(K)/B_k(K).
\]

In concrete examples:
\begin{itemize}
	\item the filled triangle has trivial \(H_1\) because its boundary loop bounds a \(2\)-simplex,
	\item the boundary of the triangle has \(H_1\cong \mathbb{Z}\) because the loop no longer bounds,
	\item the boundary of a tetrahedron has \(H_2\cong \mathbb{Z}\) because the whole surface is a nontrivial \(2\)-cycle,
	\item the torus has \(H_1\cong \mathbb{Z}^2\) and \(H_2\cong \mathbb{Z}\), reflecting its two basic loops and its fundamental surface class.
\end{itemize}

Thus simplicial homology gives a precise algebraic language for understanding the geometry of holes.

\end{document}
